\begin{register}{H}{SPI\_CMD\_REG}{\hexprefix{}0000}\label{regdesc:SPICMDREG}
\regfield{(reserved)}{7}{25}{0000000}%
\regfield{SPI\_USR}{1}{24}{{0}}%
\regfield{(reserved)}{1}{23}{0}%
\regfield{SPI\_CONF\_BITLEN}{23}{0}{{0}}%
\reglabel{Reset}\regnewline%
\begin{regdesc}\begin{reglist}
\label{fielddesc:SPICONFBITLEN}\item [SPI\_CONF\_BITLEN] 定义 SPI CONF 阶段的 SPI CLK 周期。可在 CONF 阶段配置。（读/写）
\label{fielddesc:SPIUSR}\item [SPI\_USR] 使能用户自定义命令。置位此位将触发一次 SPI 操作。操作完成后，此位清零。1：使能此功能；0：禁用此功能。CONF\_buf 不可更改该配置。（读/写）
\end{reglist}\end{regdesc}
\end{register}


\begin{register}{H}{SPI\_ADDR\_REG}{\hexprefix{}0004}\label{regdesc:SPIADDRREG}
\regfield{SPI\_USR\_ADDR\_VALUE}{32}{0}{{\hexprefix{}000000}}%
\reglabel{Reset}\regnewline%
\begin{regdesc}\begin{reglist}
\label{fielddesc:SPIUSRADDRVALUE}\item [SPI\_USR\_ADDR\_VALUE] [31:8] 位表示从机地址；[7:0]位为保留位。可在 CONF 阶段配置。（读/写）
\end{reglist}\end{regdesc}
\end{register}


\begin{register}{H}{SPI\_USER\_REG}{\hexprefix{}0018}\label{regdesc:SPIUSERREG}
\regfield{SPI\_USR\_COMMAND}{1}{31}{{1}}%
\regfield{SPI\_USR\_ADDR}{1}{30}{{0}}%
\regfield{SPI\_USR\_DUMMY}{1}{29}{{0}}%
\regfield{SPI\_USR\_MISO}{1}{28}{{0}}%
\regfield{SPI\_USR\_MOSI}{1}{27}{{0}}%
\regfield{SPI\_USR\_DUMMY\_IDLE}{1}{26}{{0}}%
\regfield{SPI\_USR\_MOSI\_HIGHPART}{1}{25}{{0}}%
\regfield{SPI\_USR\_MISO\_HIGHPART}{1}{24}{{0}}%
\regfield{(reserved)}{6}{18}{000000}%
\regfield{SPI\_USR\_HOLD\_POL}{1}{17}{{0}}%
\regfield{SPI\_SIO}{1}{16}{{0}}%
\regfield{SPI\_USR\_CONF\_NXT}{1}{15}{{0}}%
\regfield{SPI\_FWRITE\_OCT}{1}{14}{{0}}%
\regfield{SPI\_FWRITE\_QUAD}{1}{13}{{0}}%
\regfield{SPI\_FWRITE\_DUAL}{1}{12}{{0}}%
\regfield{SPI\_WR\_BYTE\_ORDER}{1}{11}{{0}}%
\regfield{SPI\_RD\_BYTE\_ORDER}{1}{10}{{0}}%
\regfield{SPI\_CK\_OUT\_EDGE}{1}{9}{{0}}%
\regfield{SPI\_RSCK\_I\_EDGE}{1}{8}{{0}}%
\regfield{SPI\_CS\_SETUP}{1}{7}{{1}}%
\regfield{SPI\_CS\_HOLD}{1}{6}{{1}}%
\regfield{SPI\_TSCK\_I\_EDGE}{1}{5}{{0}}%
\regfield{SPI\_OPI\_MODE}{1}{4}{{0}}%
\regfield{SPI\_QPI\_MODE}{1}{3}{{0}}%
\regfield{(reserved)}{2}{1}{00}%
\regfield{SPI\_DOUTDIN}{1}{0}{{0}}%
\reglabel{Reset}\regnewline%
\begin{regdesc}\begin{reglist}
\label{fielddesc:SPIDOUTDIN}\item [SPI\_DOUTDIN] 配置全双工通信。1：使能全双工通信；0：禁用全双工通信。可在 CONF 阶段配置。（读/写）
\label{fielddesc:SPIQPIMODE}\item [SPI\_QPI\_MODE] 配置 QPI 模式，适用于主机模式及从机模式。1：SPI 控制器使能 QPI 模式；0：SPI 控制器使用其它模式。可在 CONF 阶段配置。（读/写）
\label{fielddesc:SPIOPIMODE}\item [SPI\_OPI\_MODE] 配置 OPI 模式，仅用于主机模式。1：SPI 控制器使能 OPI 模式（全部为 8 位模式）。0：SPI 控制器使用其它模式。可在 CONF 阶段配置。（读/写）
\label{fielddesc:SPITSCKIEDGE}\item [SPI\_TSCK\_I\_EDGE] 从机模式下，该位用于改变 TSCK 极性。0：TSCK = SPI\_CK\_I。1：TSCK = !SPI\_CK\_I。（读/写）
\label{fielddesc:SPICSHOLD}\item [SPI\_CS\_HOLD] SPI 处于结束阶段时，SPI CS 拉低。1：使能此功能；0：禁用此功能。可在 CONF 阶段配置。（读/写）
\label{fielddesc:SPICSSETUP}\item [SPI\_CS\_SETUP] SPI 处于准备阶段时，使能 SPI CS。1：使能此功能；0：禁用此功能。可在 CONF 阶段配置。（读/写）
\label{fielddesc:SPIRSCKIEDGE}\item [SPI\_RSCK\_I\_EDGE] 从机模式下，该位用于改变 RSCK 极性。0：RSCK = !SPI\_CK\_I。1：RSCK = SPI\_CK\_I。（读/写）
\label{fielddesc:SPICKOUTEDGE}\item [SPI\_CK\_OUT\_EDGE] 与寄存器 \hyperref[regdesc:SPIDOUTMODEREG]{SPI\_DOUT\_MODE\_REG} 结合使用，配置 MOSI 信号延迟模式。可在 CONF 阶段配置。（读/写）
\label{fielddesc:SPIRDBYTEORDER}\item [SPI\_RD\_BYTE\_ORDER] 在读数据阶段，即主机输入从机输出 (MISO) 时，读取数据的字节顺序。1：选择大端序，即先读高字节；0：选择小端序，即先读低字节。可在 CONF 阶段配置。（读/写）
\label{fielddesc:SPIWRBYTEORDER}\item [SPI\_WR\_BYTE\_ORDER] 在命令阶段、地址阶段和写数据阶段，即主机输出从机输入 (MOSI) 时，选择的字节顺序。1：选择大端序，即先写高字节；0：选择小端序，即先写低字节。可在 CONF 阶段配置。（读/写）
\label{fielddesc:SPIFWRITEDUAL}\item [SPI\_FWRITE\_DUAL] 在写操作中，采用 2 位模式读取数据。可在 CONF 阶段配置。（读/写）
\label{fielddesc:SPIFWRITEQUAD}\item [SPI\_FWRITE\_QUAD] 在写操作中，采用 4 位模式读取数据。可在 CONF 阶段配置。（读/写）
\label{fielddesc:SPIFWRITEOCT}\item [SPI\_FWRITE\_OCT] 在写操作中，采用 8 位模式读取数据。可在 CONF 阶段配置。（读/写）
\label{fielddesc:SPIUSRCONFNXT}\item [SPI\_USR\_CONF\_NXT] 1：启用下次分段配置传输操作的 DMA CONF 阶段，即分段配置传输将继续进行。0：分段配置传输功能未启用，或当前一次分段配置传输结束以后，分段配置传输将不再继续进行。可在 CONF 阶段配置。（读/写）
\label{fielddesc:SPISIO}\item [SPI\_SIO] 配置三线半双工通信。MOSI 和 MISO 信号共用一个管脚。1：使能三线半双工通信；0：禁用三线半双工通信。可在 CONF 阶段配置。（读/写）

\item [见下页]
\end{reglist}\end{regdesc}
\end{register}


\addtocounter{Regfloat}{-1}
\begin{register}{H}{SPI\_USER\_REG}{\hexprefix{}0018}
\begin{regdesc}\begin{reglist}
\item [接上页]

\label{fielddesc:SPIUSRHOLDPOL}\item [SPI\_USR\_HOLD\_POL] 与保持位结合使用来设置 SPI 保持线的极性。1：保持线拉高，SPI 将处于保持状态；0：保持线拉低，SPI 将处于保持状态。可在 CONF 阶段配置。（读/写）
\label{fielddesc:SPIUSRMISOHIGHPART}\item [SPI\_USR\_MISO\_HIGHPART] 读数据阶段仅访问高位 buffer：SPI\_BUF8 \verb+~+ SPI\_BUF17。1：使能此功能；0：禁用此功能。可在 CONF 阶段配置。（读/写）
\label{fielddesc:SPIUSRMOSIHIGHPART}\item [SPI\_USR\_MOSI\_HIGHPART] 写数据阶段仅访问高位 buffer：SPI\_BUF8 \verb+~+ SPI\_BUF17。1：使能此功能；0：禁用此功能。可在 CONF 阶段配置。（读/写）
\label{fielddesc:SPIUSRDUMMYIDLE}\item [SPI\_USR\_DUMMY\_IDLE] 置位此位，则在 DUMMY 阶段禁用 SPI 时钟。可在 CONF 阶段配置。（读/写）
\label{fielddesc:SPIUSRMOSI}\item [SPI\_USR\_MOSI] 使能 SPI 操作的写数据阶段。可在 CONF 阶段配置。（读/写）
\label{fielddesc:SPIUSRMISO}\item [SPI\_USR\_MISO] 使能 SPI 操作的读数据阶段。可在 CONF 阶段配置。（读/写）
\label{fielddesc:SPIUSRDUMMY}\item [SPI\_USR\_DUMMY] 使能 SPI 操作的 DUMMY 阶段。可在 CONF 阶段配置。（读/写）
\label{fielddesc:SPIUSRADDR}\item [SPI\_USR\_ADDR] 使能 SPI 操作的地址阶段。可在 CONF 阶段配置。（读/写）
\label{fielddesc:SPIUSRCOMMAND}\item [SPI\_USR\_COMMAND] 使能 SPI 操作的命令阶段。可在 CONF 阶段配置。（读/写）
\end{reglist}\end{regdesc}
\end{register}


\begin{register}{H}{SPI\_USER1\_REG}{\hexprefix{}001C}\label{regdesc:SPIUSER1REG}
\regfield{SPI\_USR\_ADDR\_BITLEN}{5}{27}{{23}}%
\regfield{(reserved)}{19}{8}{0000000000000000000}%
\regfield{SPI\_USR\_DUMMY\_CYCLELEN}{8}{0}{{7}}%
\reglabel{Reset}\regnewline%
\begin{regdesc}\begin{reglist}
\label{fielddesc:SPIUSRDUMMYCYCLELEN}\item [SPI\_USR\_DUMMY\_CYCLELEN] 设置 DUMMY 阶段的时长，单位：spi\_clk 周期。寄存器值应为 cycle\_num（需要的时长） - 1。可在 CONF 阶段配置。（读/写）
\label{fielddesc:SPIUSRADDRBITLEN}\item [SPI\_USR\_ADDR\_BITLEN] 设置地址阶段的位长。寄存器值应为 bit\_num（需要的位长） - 1。可在 CONF 阶段配置。（读/写）
\end{reglist}\end{regdesc}
\end{register}


\begin{register}{H}{SPI\_USER2\_REG}{\hexprefix{}0020}\label{regdesc:SPIUSER2REG}
\regfield{SPI\_USR\_COMMAND\_BITLEN}{4}{28}{{7}}%
\regfield{(reserved)}{12}{16}{000000000000}%
\regfield{SPI\_USR\_COMMAND\_VALUE}{16}{0}{{0}}%
\reglabel{Reset}\regnewline%
\begin{regdesc}\begin{reglist}
\label{fielddesc:SPIUSRCOMMANDVALUE}\item [SPI\_USR\_COMMAND\_VALUE] 命令值。可在 CONF 阶段配置。（读/写）
\label{fielddesc:SPIUSRCOMMANDBITLEN}\item [SPI\_USR\_COMMAND\_BITLEN] 命令阶段的位长。寄存器值应为 bit\_num（需要的位长）- 1。可在 CONF 阶段配置。（读/写）
\end{reglist}\end{regdesc}
\end{register}


\begin{register}{H}{SPI\_MOSI\_DLEN\_REG}{\hexprefix{}0024}\label{regdesc:SPIMOSIDLENREG}
\regfield{(reserved)}{9}{23}{000000000}%
\regfield{SPI\_USR\_MOSI\_DBITLEN}{23}{0}{{\hexprefix{}0000}}%
\reglabel{Reset}\regnewline%
\begin{regdesc}\begin{reglist}
\label{fielddesc:SPIUSRMOSIDBITLEN}\item [SPI\_USR\_MOSI\_DBITLEN] 写数据的位长。寄存器值应为 bit\_num（需要的位长）- 1。可在 CONF 阶段配置。（读/写）
\end{reglist}\end{regdesc}
\end{register}


\begin{register}{H}{SPI\_MISO\_DLEN\_REG}{\hexprefix{}0028}\label{regdesc:SPIMISODLENREG}
\regfield{(reserved)}{9}{23}{000000000}%
\regfield{SPI\_USR\_MISO\_DBITLEN}{23}{0}{{\hexprefix{}0000}}%
\reglabel{Reset}\regnewline%
\begin{regdesc}\begin{reglist}
\label{fielddesc:SPIUSRMISODBITLEN}\item [SPI\_USR\_MISO\_DBITLEN] 读数据的位长。寄存器值应为 bit\_num（需要的位长）- 1。可在 CONF 阶段配置。（读/写）
\end{reglist}\end{regdesc}
\end{register}


\begin{register}{H}{SPI\_CTRL\_REG}{\hexprefix{}0008}\label{regdesc:SPICTRLREG}
\regfield{(reserved)}{5}{27}{00000}%
\regfield{SPI\_WR\_BIT\_ORDER}{1}{26}{{0}}%
\regfield{SPI\_RD\_BIT\_ORDER}{1}{25}{{0}}%
\regfield{(reserved)}{3}{22}{000}%
\regfield{SPI\_WP\_REG}{1}{21}{{1}}%
\regfield{(reserved)}{1}{20}{0}%
\regfield{SPI\_D\_POL}{1}{19}{{1}}%
\regfield{SPI\_Q\_POL}{1}{18}{{1}}%
\regfield{(reserved)}{1}{17}{0}%
\regfield{SPI\_FREAD\_OCT}{1}{16}{{0}}%
\regfield{SPI\_FREAD\_QUAD}{1}{15}{{0}}%
\regfield{SPI\_FREAD\_DUAL}{1}{14}{{0}}%
\regfield{(reserved)}{3}{11}{000}%
\regfield{SPI\_FCMD\_OCT}{1}{10}{{0}}%
\regfield{SPI\_FCMD\_QUAD}{1}{9}{{0}}%
\regfield{SPI\_FCMD\_DUAL}{1}{8}{{0}}%
\regfield{SPI\_FADDR\_OCT}{1}{7}{{0}}%
\regfield{SPI\_FADDR\_QUAD}{1}{6}{{0}}%
\regfield{SPI\_FADDR\_DUAL}{1}{5}{{0}}%
\regfield{(reserved)}{1}{4}{0}%
\regfield{SPI\_DUMMY\_OUT}{1}{3}{{0}}%
\regfield{(reserved)}{3}{0}{000}%
\reglabel{Reset}\regnewline%
\begin{regdesc}\begin{reglist}
\label{fielddesc:SPIDUMMYOUT}\item [SPI\_DUMMY\_OUT] 在 DUMMY 阶段，SPI 信号电平由 SPI 控制器输出。可在 CONF 阶段配置。（读/写）
\label{fielddesc:SPIFADDRDUAL}\item [SPI\_FADDR\_DUAL] 在地址阶段，采用 2 位模式。1：使能 2 位模式；0：禁用 2 位模式。可在 CONF 阶段配置。（读/写）
\label{fielddesc:SPIFADDRQUAD}\item [SPI\_FADDR\_QUAD] 在地址阶段，采用 4 位模式。1：使能 4 位模式；0：禁用 4 位模式。可在 CONF 阶段配置。（读/写）
\label{fielddesc:SPIFADDROCT}\item [SPI\_FADDR\_OCT] 在地址阶段，采用 8 位模式。1：使能 8 位模式；0：禁用 8 位模式。可在 CONF 阶段配置。（读/写）
\label{fielddesc:SPIFCMDDUAL}\item [SPI\_FCMD\_DUAL] 在命令阶段，采用 2 位模式。1：使能 2 位模式；0：禁用 2 位模式。可在 CONF 阶段配置。（读/写）
\label{fielddesc:SPIFCMDQUAD}\item [SPI\_FCMD\_QUAD] 在命令阶段，采用 4 位模式。1：使能 4 位模式；0：禁用 4 位模式。可在 CONF 阶段配置。（读/写）
\label{fielddesc:SPIFCMDOCT}\item [SPI\_FCMD\_OCT] 在命令阶段，采用 8 位模式。1：使能 8 位模式；0：禁用 8 位模式。可在 CONF 阶段配置。（读/写）
\label{fielddesc:SPIFREADDUAL}\item [SPI\_FREAD\_DUAL] 在读操作中，采用 2 位模式读取数据。1：使能 2 位模式；0：禁用 2 位模式。可在 CONF 阶段配置。（读/写）
\label{fielddesc:SPIFREADQUAD}\item [SPI\_FREAD\_QUAD] 在读操作中，采用 4 位模式读取数据。1：使能 4 位模式；0：禁用 4 位模式。可在 CONF 阶段配置。（读/写）
\label{fielddesc:SPIFREADOCT}\item [SPI\_FREAD\_OCT] 在读操作中，采用 8 位模式读取数据。1：使能 8 位模式；0：禁用 8 位模式。可在 CONF 阶段配置。（读/写）
\label{fielddesc:SPIQPOL}\item [SPI\_Q\_POL] 配置 MISO 线的极性，1：高；0：低。可在 CONF 阶段配置。（读/写）
\label{fielddesc:SPIDPOL}\item [SPI\_D\_POL] 配置 MOSI 线的极性，1：高；0：低。可在 CONF 阶段配置。（读/写）
\label{fielddesc:SPIWPREG}\item [SPI\_WP\_REG] SPI 处于空闲状态时，输出写保护信号。1：输出高电平；0：输出低电平。可在 CONF 阶段配置。（读/写）
\label{fielddesc:SPIRDBITORDER}\item [SPI\_RD\_BIT\_ORDER] 在读数据阶段，即主机输入从机输出 (MISO) 时，读取数据的位顺序。1：先读最低有效位 (LSB)；0：先读最高有效位 (MSB)。可在 CONF 阶段配置。（读/写）
\label{fielddesc:SPIWRBITORDER}\item [SPI\_WR\_BIT\_ORDER] 在命令阶段、地址阶段和写数据阶段，即主机输出从机输入 (MOSI) 时，写数据的位顺序。1：先写最低有效为 (LSB)；0：先写最高有效位 (MSB)。可在 CONF 阶段配置。（读/写）
\end{reglist}\end{regdesc}
\end{register}


\begin{register}{H}{SPI\_CTRL1\_REG}{\hexprefix{}000C}\label{regdesc:SPICTRL1REG}
\regfield{(reserved)}{12}{20}{000000000000}%
\regfield{SPI\_CS\_HOLD\_DELAY}{6}{14}{{\hexprefix{}1}}%
\regfield{(reserved)}{9}{5}{000000000}%
\regfield{SPI\_W16\_17\_WR\_ENA}{1}{4}{{1}}%
\regfield{SPI\_RSCK\_DATA\_OUT}{1}{3}{{0}}%
\regfield{SPI\_CLK\_MODE\_13}{1}{2}{{0}}%
\regfield{SPI\_CLK\_MODE}{2}{0}{{\hexprefix{}0}}%
\reglabel{Reset}\regnewline%
\begin{regdesc}\begin{reglist}
\label{fielddesc:SPICLKMODE}\item [SPI\_CLK\_MODE] SPI 时钟模式控制位。0：CS 拉高时，SPI 时钟关闭；1：CS 拉高后，SPI 延迟一个时钟周期；2：CS 拉高后，SPI 延迟两个时钟周期；3：SPI 时钟一直处于工作状态。可在 CONF 阶段配置。（读/写）
\label{fielddesc:SPICLKMODE13}\item [SPI\_CLK\_MODE\_13] {CPOL, CPHA} 1：支持 SPI CLK 模式 1 和模式 3，第一边沿输出数据 B[0]/B[7]。1：支持 SPI CLK 模式 0 和 模式 2，第一边沿输出数据 B[1]/B[6]。（读/写）

\label{fielddesc:SPIRSCKDATAOUT}\item [SPI\_RSCK\_DATA\_OUT] 当 TSCK 与 RSCK 相同时，将节省半个周期。1：在 RSCK 上升沿输出数据；0：在 TSCK 上升沿输出数据。（读/写）

\label{fielddesc:SPIW1617WRENA}\item [SPI\_W16\_17\_WR\_ENA] 1：可以写 SPI\_BUF16 \verb+~+ SPI\_BUF17；0：不可以写 SPI\_BUF16 \verb+~+ SPI\_BUF17。可在 CONF 阶段配置。（读/写）
\label{fielddesc:SPICSHOLDDELAY}\item [SPI\_CS\_HOLD\_DELAY] SPI CS 信号被延时的 SPI 时钟周期数。可在 CONF 阶段配置。（读/写）
\end{reglist}\end{regdesc}
\end{register}


\begin{register}{H}{SPI\_CTRL2\_REG}{\hexprefix{}0010}\label{regdesc:SPICTRL2REG}
\regfield{(reserved)}{1}{31}{0}%
\regfield{SPI\_CS\_DELAY\_NUM}{2}{29}{{\hexprefix{}0}}%
\regfield{SPI\_CS\_DELAY\_MODE}{3}{26}{{\hexprefix{}0}}%
\regfield{SPI\_CS\_HOLD\_TIME}{13}{13}{{\hexprefix{}01}}%
\regfield{SPI\_CS\_SETUP\_TIME}{13}{0}{{\hexprefix{}00}}%
\reglabel{Reset}\regnewline%
\begin{regdesc}\begin{reglist}
\label{fielddesc:SPICSSETUPTIME}\item [SPI\_CS\_SETUP\_TIME] 准备阶段所用的周期数 + 1，时钟为 SPI 时钟。此位与 \hyperref[fielddesc:SPICSSETUP]{SPI\_CS\_SETUP} 搭配使用。可在 CONF 阶段配置。（读/写）
\label{fielddesc:SPICSHOLDTIME}\item [SPI\_CS\_HOLD\_TIME] CS 的延迟周期数，时钟为 SPI 时钟。此位与 \hyperref[fielddesc:SPICSSETUP]{SPI\_CS\_HOLD} 搭配使用。可在 CONF 阶段配置。（读/写）
\label{fielddesc:SPICSDELAYMODE}\item [SPI\_CS\_DELAY\_MODE] SPI CS 信号被延迟，延迟单位：SPI CLK。0：延迟 0 个周期；1：如果置位 \hyperref[fielddesc:SPICKOUTEDGE]{SPI\_CK\_OUT\_EDGE} 或 \hyperref[fielddesc:SPICKIDLEEDGE]{SPI\_CK\_IDLE\_EDGE}，则延迟半个周期，否则延迟一个周期；2：如果置位 \hyperref[fielddesc:SPICKOUTEDGE]{SPI\_CK\_OUT\_EDGE} 或 \hyperref[fielddesc:SPICKIDLEEDGE]{SPI\_CK\_IDLE\_EDGE}，则延迟一个周期，否则延迟半周期；3：延迟一个周期。可在 CONF 阶段配置。（读/写）
\label{fielddesc:SPICSDELAYNUM}\item [SPI\_CS\_DELAY\_NUM] SPI CS 信号被延迟，延迟单位：系统时钟周期。可在 CONF 阶段配置。（读/写）
\end{reglist}\end{regdesc}
\end{register}


\begin{register}{H}{SPI\_CLOCK\_REG}{\hexprefix{}0014}\label{regdesc:SPICLOCKREG}
\regfield{SPI\_CLK\_EQU\_SYSCLK}{1}{31}{{1}}%
\regfield{SPI\_CLKDIV\_PRE}{13}{18}{{0}}%
\regfield{SPI\_CLKCNT\_N}{6}{12}{{\hexprefix{}3}}%
\regfield{SPI\_CLKCNT\_H}{6}{6}{{\hexprefix{}1}}%
\regfield{SPI\_CLKCNT\_L}{6}{0}{{\hexprefix{}3}}%
\reglabel{Reset}\regnewline%
\begin{regdesc}\begin{reglist}
\label{fielddesc:SPICLKCNTL}\item [SPI\_CLKCNT\_L] 主机模式下，必须与 \hyperref[fielddesc:SPICLKCNTN]{SPI\_CLKCNT\_N} 相等。从机模式下，必须为 0。可在 CONF 阶段配置。（读/写）
\label{fielddesc:SPICLKCNTH}\item [SPI\_CLKCNT\_H] 主机模式下，必须为 (\hyperref[fielddesc:SPICLKCNTN]{SPI\_CLKCNT\_N} + 1)/2-1 的向下取整值。从机模式下，必须为 0。可在 CONF 阶段配置。（读/写）
\label{fielddesc:SPICLKCNTN}\item [SPI\_CLKCNT\_N] 主机模式下，SPI CLK 的分频系数。因此 SPI CLK 频率为 $f_{\textrm{apb}}$/(\hyperref[fielddesc:SPICLKDIVPRE]{SPI\_CLKDIV\_PRE} + 1)/(\hyperref[fielddesc:SPICLKCNTN]{SPI\_CLKCNT\_N} + 1)。可在 CONF 阶段配置。（读/写）
\label{fielddesc:SPICLKDIVPRE}\item [SPI\_CLKDIV\_PRE] 主机模式下，SPI CLK 的预分频系数。可在 CONF 阶段配置。（读/写）
\label{fielddesc:SPICLKEQUSYSCLK}\item [SPI\_CLK\_EQU\_SYSCLK] 在主机模式下，1：SPI CLK 与 APB CLK 频率相同；0：SPI CLK 从 APB CLK 分频产生。可在 CONF 阶段配置。（读/写）
\end{reglist}\end{regdesc}
\end{register}


\begin{register}{H}{SPI\_MISC\_REG}{\hexprefix{}002C}\label{regdesc:SPIMISCREG}
\regfield{SPI\_QUAD\_DIN\_PIN\_SWAP}{1}{31}{{0}}%
\regfield{SPI\_CS\_KEEP\_ACTIVE}{1}{30}{{0}}%
\regfield{SPI\_CK\_IDLE\_EDGE}{1}{29}{{0}}%
\regfield{(reserved)}{2}{27}{00}%
\regfield{SPI\_CD\_IDLE\_EDGE}{1}{26}{{0}}%
\regfield{SPI\_CD\_CMD\_SET}{1}{25}{{0}}%
\regfield{SPI\_DQS\_IDLE\_EDGE}{1}{24}{{0}}%
\regfield{SPI\_SLAVE\_CS\_POL}{1}{23}{{0}}%
\regfield{SPI\_CD\_ADDR\_SET}{1}{22}{{0}}%
\regfield{SPI\_CD\_DUMMY\_SET}{1}{21}{{0}}%
\regfield{SPI\_CD\_DATA\_SET}{1}{20}{{0}}%
\regfield{SPI\_CMD\_DTR\_EN}{1}{19}{{0}}%
\regfield{SPI\_ADDR\_DTR\_EN}{1}{18}{{0}}%
\regfield{SPI\_DATA\_DTR\_EN}{1}{17}{{0}}%
\regfield{SPI\_CLK\_DATA\_DTR\_EN}{1}{16}{{0}}%
\regfield{(reserved)}{3}{13}{000}%
\regfield{SPI\_MASTER\_CS\_POL}{6}{7}{{0}}%
\regfield{SPI\_CK\_DIS}{1}{6}{{0}}%
\regfield{SPI\_CS5\_DIS}{1}{5}{{1}}%
\regfield{SPI\_CS4\_DIS}{1}{4}{{1}}%
\regfield{SPI\_CS3\_DIS}{1}{3}{{1}}%
\regfield{SPI\_CS2\_DIS}{1}{2}{{1}}%
\regfield{SPI\_CS1\_DIS}{1}{1}{{1}}%
\regfield{SPI\_CS0\_DIS}{1}{0}{{0}}%
\reglabel{Reset}\regnewline%
\begin{regdesc}\begin{reglist}
\label{fielddesc:SPICS0DIS}\item [SPI\_CS0\_DIS] SPI CS0 管脚使能；1：禁用 CS0；0：SPI CS0 信号引入/引出 CS0 管脚。可在 CONF 阶段配置。（读/写）
\label{fielddesc:SPICS1DIS}\item [SPI\_CS1\_DIS] SPI CS1 管脚使能；1：禁用 CS1；0：SPI CS1 信号引入/引出 CS1 管脚。可在 CONF 阶段配置。（读/写）
\label{fielddesc:SPICS2DIS}\item [SPI\_CS2\_DIS] SPI CS2 管脚使能；1：禁用 CS2；0：SPI CS2 信号引入/引出 CS2 管脚。可在 CONF 阶段配置。（读/写）
\label{fielddesc:SPICS3DIS}\item [SPI\_CS3\_DIS] SPI CS3 管脚使能；1：禁用 CS3；0：SPI CS3 信号引入/引出 CS3 管脚。可在 CONF 阶段配置。（读/写）
\label{fielddesc:SPICS4DIS}\item [SPI\_CS4\_DIS] SPI CS4 管脚使能；1：禁用 CS4；0：SPI CS4 信号引入/引出 CS4 管脚。可在 CONF 阶段配置。（读/写）
\label{fielddesc:SPICS5DIS}\item [SPI\_CS5\_DIS] SPI CS5 管脚使能；1：禁用 CS5；0：SPI CS5 信号引入/引出 CS5 管脚。可在 CONF 阶段配置。（读/写）
\label{fielddesc:SPICKDIS}\item [SPI\_CK\_DIS] 1：停止 SPI CLK 输出；0：使能 SPI CLK 输出。可在 CONF 阶段配置。（读/写）
\label{fielddesc:SPIMASTERCSPOL}\item [SPI\_MASTER\_CS\_POL] 在主机模式下，SPI CS 线的极性，该值等于 SPI\_CS\^{} SPI\_MASTER\_CS\_POL。可在 CONF 阶段配置。（读/写）

\label{fielddesc:SPICLKDATADTREN}\item [SPI\_CLK\_DATA\_DTR\_EN] 1：SPI CLK、SPI 数据和 SPI DQS 采用 SPI 主机 DTR 模式。0：仅 SPI DQS 采用 SPI 主机 DTR 模式。此位结合 SPI\_DATA\_DTR\_EN、SPI\_ADDR\_DTR\_EN、SPI\_CMD\_DTR\_EN 一起使用。（读/写）
\label{fielddesc:SPIDATADTREN}\item [SPI\_DATA\_DTR\_EN] 1：SPI CLK 及 SPI DOUT 和 SPI DIN 阶段的数据使用 DTR 模式，包括主机 1/2/4/8 位模式。0：SPI CLK 及 SPI DOUT 和 SPI DIN 阶段的数据使用 STR 模式。可在 CONF 阶段配置。（读/写）
\label{fielddesc:SPIADDRDTREN}\item [SPI\_ADDR\_DTR\_EN] 1：SPI CLK 和 SPI\_SEND\_ADDR 阶段的数据使用 DTR 模式，包括主机 1/2/4/8 位模式。0：SPI CLK 和 SPI\_SEND\_ADDR 阶段的数据使用 STR 模式。可在 CONF 阶段配置。（读/写）
\label{fielddesc:SPICMDDTREN}\item [SPI\_CMD\_DTR\_EN] 1：SPI CLK 和 SPI\_SEND\_CMD 阶段的数据使用 DTR 模式，包括主机 1/2/4/8 位模式。0：SPI CLK 和 SPI\_SEND\_CMD 阶段的数据使用 STR 模式。可在 CONF 阶段配置。（读/写）
\label{fielddesc:SPICDDATASET}\item [SPI\_CD\_DATA\_SET] 1：\hyperref[fielddesc:SPIST]{SPI\_ST}[3:0] 处于 SPI\_DOUT 或 SPI\_DIN 状态时，SPI\_CD = !SPI\_CD\_IDLE\_EDGE。0：SPI\_CD = SPI\_CD\_IDLE\_EDGE。可在 CONF 阶段配置。（读/写）
\item [见下页]
\end{reglist}\end{regdesc}
\end{register}


\addtocounter{Regfloat}{-1}
\begin{register}{H}{SPI\_MISC\_REG}{\hexprefix{}002C}
\begin{regdesc}\begin{reglist}
\item [接上页]
\label{fielddesc:SPICDDUMMYSET}\item [SPI\_CD\_DUMMY\_SET] 1: SPI\_ST[3:0] 处于 SPI\_DUMMY 状态时，SPI\_CD = !SPI\_CD\_IDLE\_EDGE。0：SPI\_CD = SPI\_CD\_IDLE\_EDGE。可在 CONF 阶段配置。（读/写）
\label{fielddesc:SPICDADDRSET}\item [SPI\_CD\_ADDR\_SET] 1: SPI\_ST[3:0] 处于 SPI\_SEND\_ADDR 状态时，SPI\_CD = !SPI\_CD\_IDLE\_EDGE。0：SPI\_CD = SPI\_CD\_IDLE\_EDGE。可在 CONF 阶段配置。（读/写）
\label{fielddesc:SPISLAVECSPOL}\item [SPI\_SLAVE\_CS\_POL] SPI 从机输入 CS 极性选择。1：反相；0：保持不变。可在 CONF 阶段配置。（读/写）
\label{fielddesc:SPIDQSIDLEEDGE}\item [SPI\_DQS\_IDLE\_EDGE] 默认值为 SPI\_DQS。可在 CONF 阶段配置。（读/写）
\label{fielddesc:SPICDCMDSET}\item [SPI\_CD\_CMD\_SET] 1: SPI\_ST[3:0] 处于 SPI\_SEND\_CMD 状态时，SPI\_CD = !SPI\_CD\_IDLE\_EDGE。0：SPI\_CD = SPI\_CD\_IDLE\_EDGE。可在 CONF 阶段配置。（读/写）
\label{fielddesc:SPICDIDLEEDGE}\item [SPI\_CD\_IDLE\_EDGE] 默认值为 SPI\_CD。可在 CONF 阶段配置。（读/写）
\label{fielddesc:SPICKIDLEEDGE}\item [SPI\_CK\_IDLE\_EDGE] 1：空闲状态时，SPI CLK 线拉高；0：空闲状态时，SPI CLK 线拉低。可在 CONF 阶段配置。（读/写）
\label{fielddesc:SPICSKEEPACTIVE}\item [SPI\_CS\_KEEP\_ACTIVE] 置位此位，则 SPI CS 线拉低。可在 CONF 阶段配置。（读/写）
\label{fielddesc:SPIQUADDINPINSWAP}\item [SPI\_QUAD\_DIN\_PIN\_SWAP] 1：使能 SPI Quad 输入交换；0：禁止 SPI Quad 输入交换。可在 CONF 阶段配置。（读/写）
\end{reglist}\end{regdesc}
\end{register}


\begin{register}{H}{SPI\_FSM\_REG}{\hexprefix{}0044}\label{regdesc:SPIFSMREG}
\regfield{SPI\_MST\_DMA\_RD\_BYTELEN}{20}{12}{{\hexprefix{}000}}%
\regfield{(reserved)}{8}{4}{00000000}%
\regfield{SPI\_ST}{4}{0}{{0}}%
\reglabel{Reset}\regnewline%
\begin{regdesc}\begin{reglist}
\label{fielddesc:SPIST}\item [SPI\_ST] SPI 状态机的状态。0：空闲状态；1：准备状态；2：发送命令状态；3：发送数据状态；4：读数据状态；5：写数据状态；6：等待状态；7：完成状态。 （只读）
\label{fielddesc:SPIMSTDMARDBYTELEN}\item [SPI\_MST\_DMA\_RD\_BYTELEN] 在非分段配置传输或分段配置传输模式下定义主机 DMA 读取字节的长度。SPI\_RX\_EOF\_EN 为 0 时无效。可在 CONF 阶段配置。（读/写）
\end{reglist}\end{regdesc}
\end{register}


\begin{register}{H}{SPI\_HOLD\_REG}{\hexprefix{}0048}\label{regdesc:SPIHOLDREG}
\regfield{(reserved)}{24}{8}{000000000000000000000000}%
\regfield{SPI\_DMA\_SEG\_TRANS\_DONE}{1}{7}{{0}}%
\regfield{SPI\_HOLD\_OUT\_TIME}{3}{4}{{0}}%
\regfield{SPI\_HOLD\_OUT\_EN}{1}{3}{{0}}%
\regfield{SPI\_HOLD\_VAL\_REG}{1}{2}{{0}}%
\regfield{(reserved)}{2}{0}{00}%
\reglabel{Reset}\regnewline%
\begin{regdesc}\begin{reglist}
\label{fielddesc:SPIHOLDVALREG}\item [SPI\_HOLD\_VAL\_REG] SPI 保持线的输出值，应与 SPI\_HOLD\_OUT\_EN 一同使用。可在 CONF 阶段配置。（读/写）
\label{fielddesc:SPIHOLDOUTEN}\item [SPI\_HOLD\_OUT\_EN] 启用将 SPI 保持值输出至 SPI\_HOLD\_EN。与 SPI\_EXT\_HOLD\_EN 和其它 USR 保持信号一起用于保持状态机。可在 CONF 阶段配置。（读/写）
\label{fielddesc:SPIHOLDOUTTIME}\item [SPI\_HOLD\_OUT\_TIME] SPI\_HOLD\_OUT\_EN 使能的情况下，设置输出 SPI 保持信号的保持周期。可在 CONF 阶段配置。（读/写）
\label{fielddesc:SPIDMASEGTRANSDONE}\item [SPI\_DMA\_SEG\_TRANS\_DONE] 1：SPI 主机 DMA 全双工/半双工分段配置传输结束，或从机半双工连续传输结束。数据已转存至相应内存。0：主机分段配置传输/从机连续传输未结束或未发生。CONF\_buf 不可更改该设置。（读/写）
\end{reglist}\end{regdesc}
\end{register}


\begin{register}{H}{SPI\_SLAVE\_REG}{\hexprefix{}0030}\label{regdesc:SPISLAVEREG}
\regfield{SPI\_SOFT\_RESET}{1}{31}{{0}}%
\regfield{SPI\_SLAVE\_MODE}{1}{30}{{0}}%
\regfield{SPI\_TRANS\_DONE\_AUTO\_CLR\_EN}{1}{29}{{0}}%
\regfield{(reserved)}{2}{27}{00}%
\regfield{SPI\_TRANS\_CNT}{4}{23}{{0}}%
\regfield{(reserved)}{11}{12}{00000000000}%
\regfield{SPI\_SEG\_MAGIC\_ERR\_INT\_EN}{1}{11}{{0}}%
\regfield{SPI\_INT\_DMA\_SEG\_TRANS\_EN}{1}{10}{{0}}%
\regfield{SPI\_INT\_TRANS\_DONE\_EN}{1}{9}{{1}}%
\regfield{SPI\_INT\_WR\_DMA\_DONE\_EN}{1}{8}{{0}}%
\regfield{SPI\_INT\_RD\_DMA\_DONE\_EN}{1}{7}{{0}}%
\regfield{SPI\_INT\_WR\_BUF\_DONE\_EN}{1}{6}{{0}}%
\regfield{SPI\_INT\_RD\_BUF\_DONE\_EN}{1}{5}{{0}}%
\regfield{SPI\_TRANS\_DONE}{1}{4}{{0}}%
\regfield{(reserved)}{4}{0}{0000}%
\reglabel{Reset}\regnewline%
\begin{regdesc}\begin{reglist}
\label{fielddesc:SPITRANSDONE}\item [SPI\_TRANS\_DONE] 主机模式和从机模式下，操作完成即触发此中断。该位为中断原始位。CONF\_buf 不可更改该设置。（读/写）
\label{fielddesc:SPIINTRDBUFDONEEN}\item [SPI\_INT\_RD\_BUF\_DONE\_EN] SPI\_SLV\_RD\_BUF\_DONE 中断使能位。1：使能此中断；0：禁用此中断。可在 CONF 阶段配置。（读/写）
\label{fielddesc:SPIINTWRBUFDONEEN}\item [SPI\_INT\_WR\_BUF\_DONE\_EN] SPI\_SLV\_WR\_BUF\_DONE 中断使能位。1：使能此中断；0：禁用此中断。可在 CONF 阶段配置。（读/写）
\label{fielddesc:SPIINTRDDMADONEEN}\item [SPI\_INT\_RD\_DMA\_DONE\_EN] SPI\_SLV\_RD\_DMA\_DONE 中断使能位。1：使能此中断；0：禁用此中断。可在 CONF 阶段配置。（读/写）
\label{fielddesc:SPIINTWRDMADONEEN}\item [SPI\_INT\_WR\_DMA\_DONE\_EN] SPI\_SLV\_WR\_DMA\_DONE 中断使能位。1：使能此中断；0：禁用此中断。可在 CONF 阶段配置。（读/写）
\label{fielddesc:SPIINTTRANSDONEEN}\item [SPI\_INT\_TRANS\_DONE\_EN] SPI\_TRANS\_DONE 中断使能位。1：使能此中断；0：禁用此中断。可在 CONF 阶段配置。（读/写）
\label{fielddesc:SPIINTDMASEGTRANSEN}\item [SPI\_INT\_DMA\_SEG\_TRANS\_EN] SPI\_DMA\_SEG\_TRANS\_DONE 中断使能位。1：使能此中断；0：禁用此中断。可在 CONF 阶段配置。（读/写）
\label{fielddesc:SPISEGMAGICERRINTEN}\item [SPI\_SEG\_MAGIC\_ERR\_INT\_EN] 1：使能 Magic 值错误中断。0：使用其它中断。可在 CONF 阶段配置。（读/写）
\label{fielddesc:SPITRANSCNT}\item [SPI\_TRANS\_CNT] 操作计数器，同时适用于主机模式和从机模式。 （只读）
\label{fielddesc:SPITRANSDONEAUTOCLREN}\item [SPI\_TRANS\_DONE\_AUTO\_CLR\_EN] 使能 SPI\_TRANS\_DONE 自动清零模式，可在 SPI\_TRANS\_DONE 上升沿后 3 个 APB 时钟清除 SPI\_TRANS\_DONE。0：禁用此功能。1：使能此功能。可在 CONF 阶段配置。（读/写）
\label{fielddesc:SPISLAVEMODE}\item [SPI\_SLAVE\_MODE] 设置 SPI 工作模式。1：从机模式；0：主机模式。（读/写）
\label{fielddesc:SPISOFTRESET}\item [SPI\_SOFT\_RESET] 使能软件置位，置位 SPI 时钟线、CS 线和数据线。可在 CONF 阶段配置。（读/写）
\end{reglist}\end{regdesc}
\end{register}


\begin{register}{H}{SPI\_SLAVE1\_REG}{\hexprefix{}0034}\label{regdesc:SPISLAVE1REG}
\regfield{SPI\_SLV\_LAST\_ADDR}{8}{24}{{0}}%
\regfield{SPI\_SLV\_LAST\_COMMAND}{8}{16}{{0}}%
\regfield{SPI\_SLV\_WR\_DMA\_DONE}{1}{15}{{0}}%
\regfield{SPI\_SLV\_CMD\_ERR}{1}{14}{{0}}%
\regfield{SPI\_SLV\_ADDR\_ERR}{1}{13}{{0}}%
\regfield{SPI\_SLV\_NO\_QPI\_EN}{1}{12}{{0}}%
\regfield{SPI\_SLV\_CMD\_ERR\_CLR}{1}{11}{{0}}%
\regfield{SPI\_SLV\_ADDR\_ERR\_CLR}{1}{10}{{0}}%
\regfield{(reserved)}{10}{0}{0000000000}%
\reglabel{Reset}\regnewline%
\begin{regdesc}\begin{reglist}
\label{fielddesc:SPISLVADDRERRCLR}\item [SPI\_SLV\_ADDR\_ERR\_CLR] 1：清除 SPI\_SLV\_ADDR\_ERR。0：无效。CONF\_buf 可以更改该设置。（读/写）
\label{fielddesc:SPISLVCMDERRCLR}\item [SPI\_SLV\_CMD\_ERR\_CLR] 1：清除 SPI\_SLV\_CMD\_ERR。0：无效。CONF\_buf 可以更改该设置。（读/写）
\label{fielddesc:SPISLVNOQPIEN}\item [SPI\_SLV\_NO\_QPI\_EN] 1：不支持从机 QPI 模式。0：支持从机 QPI 模式。（读/写）
\label{fielddesc:SPISLVADDRERR}\item [SPI\_SLV\_ADDR\_ERR] 1：SPI 从机模式不支持上次 SPI 传输的地址值。0：支持 SPI 传输的地址值或未收到地址值。 （只读）
\label{fielddesc:SPISLVCMDERR}\item [SPI\_SLV\_CMD\_ERR] 1：SPI 从机模式不支持上次 SPI 传输的命令值。0：支持命令值或未收到命令值。 （只读）
\label{fielddesc:SPISLVWRDMADONE}\item [SPI\_SLV\_WR\_DMA\_DONE] 从机模式下，完成 DMA 写操作将触发中断，该位为中断原始位。CONF\_buf 不可更改该设置。（读/写）
\label{fielddesc:SPISLVLASTCOMMAND}\item [SPI\_SLV\_LAST\_COMMAND] 从机模式下，锁存上次传输的命令值。（读/写）
\label{fielddesc:SPISLVLASTADDR}\item [SPI\_SLV\_LAST\_ADDR] 从机模式下，锁存上次传输的地址值。（读/写）
\end{reglist}\end{regdesc}
\end{register}


\begin{register}{H}{SPI\_SLV\_WRBUF\_DLEN\_REG}{\hexprefix{}0038}\label{regdesc:SPISLVWRBUFDLENREG}
\regfield{SPI\_CONF\_BASE\_BITLEN}{7}{25}{{108}}%
\regfield{SPI\_SLV\_WR\_BUF\_DONE}{1}{24}{{0}}%
\regfield{(reserved)}{24}{0}{000000000000000000000000}%
\reglabel{Reset}\regnewline%
\begin{regdesc}\begin{reglist}
\label{fielddesc:SPISLVWRBUFDONE}\item [SPI\_SLV\_WR\_BUF\_DONE] 从机模式下，DMA 写 buffer 操作完成将触发中断，该位为中断原始位。CONF\_buf 不可更改该设置。（读/写）
\label{fielddesc:SPICONFBASEBITLEN}\item [SPI\_CONF\_BASE\_BITLEN] CONF 阶段的基本周期，时钟：SPI CLK。如果使能了 \hyperref[fielddesc:SPIUSRCONF]{SPI\_USR\_CONF}，则 CONF 阶段的周期为 SPI\_CONF\_BASE\_BITLEN[6:0] + \hyperref[fielddesc:SPICONFBITLEN]{SPI\_CONF\_BITLEN}[23:0]。（读/写）
\end{reglist}\end{regdesc}
\end{register}


\begin{register}{H}{SPI\_SLV\_RDBUF\_DLEN\_REG}{\hexprefix{}003C}\label{regdesc:SPISLVRDBUFDLENREG}
\regfield{(reserved)}{6}{26}{000000}%
\regfield{SPI\_SEG\_MAGIC\_ERR}{1}{25}{{0}}%
\regfield{SPI\_SLV\_RD\_BUF\_DONE}{1}{24}{{0}}%
\regfield{(reserved)}{4}{20}{0000}%
\regfield{SPI\_SLV\_DMA\_RD\_BYTELEN}{20}{0}{{\hexprefix{}000}}%
\reglabel{Reset}\regnewline%
\begin{regdesc}\begin{reglist}
\label{fielddesc:SPISLVDMARDBYTELEN}\item [SPI\_SLV\_DMA\_RD\_BYTELEN] 从机模式下，读操作长度，单位：字节。寄存器值为 byte\_num - 1。（读/写）
\label{fielddesc:SPISLVRDBUFDONE}\item [SPI\_SLV\_RD\_BUF\_DONE] 从机模式下，DMA 读 buffer 操作完成将触发中断，该位为中断原始位。CONF\_buf 不可更改该设置。（读/写）
\label{fielddesc:SPISEGMAGICERR}\item [SPI\_SEG\_MAGIC\_ERR] 1：主机 DMA 分段配置传输模式下，CONF buffer 中最近的 Magic 值有误。0：其它模式。（读/写）
\end{reglist}\end{regdesc}
\end{register}


\begin{register}{H}{SPI\_SLV\_RD\_BYTE\_REG}{\hexprefix{}0040}\label{regdesc:SPISLVRDBYTEREG}
\regfield{SPI\_USR\_CONF}{1}{31}{{0}}%
\regfield{SPI\_SLV\_RD\_DMA\_DONE}{1}{30}{{0}}%
\regfield{(reserved)}{2}{28}{00}%
\regfield{SPI\_DMA\_SEG\_MAGIC\_VALUE}{4}{24}{{10}}%
\regfield{SPI\_SLV\_WRBUF\_BYTELEN\_EN}{1}{23}{{0}}%
\regfield{SPI\_SLV\_RDBUF\_BYTELEN\_EN}{1}{22}{{0}}%
\regfield{SPI\_SLV\_WRDMA\_BYTELEN\_EN}{1}{21}{{0}}%
\regfield{SPI\_SLV\_RDDMA\_BYTELEN\_EN}{1}{20}{{0}}%
\regfield{SPI\_SLV\_DATA\_BYTELEN}{20}{0}{{0}}%
\reglabel{Reset}\regnewline%
\begin{regdesc}\begin{reglist}
\label{fielddesc:SPISLVDATABYTELEN}\item [SPI\_SLV\_DATA\_BYTELEN] 从机模式下，上次 SPI 传输的全双工或半双工数据字节长度。半双工模式下，该值由 SPI\_SLV\_RDDMA\_BYTELEN\_EN、SPI\_SLV\_WRDMA\_BYTELEN\_EN、SPI\_SLV\_RDBUF\_BYTELEN\_EN、SPI\_SLV\_WRBUF\_BYTELEN\_EN 共同控制。（读/写）
\label{fielddesc:SPISLVRDDMABYTELENEN}\item [SPI\_SLV\_RDDMA\_BYTELEN\_EN] 1：DMA 控制模式下 (Rd\_DMA)，SPI\_SLV\_DATA\_BYTELEN 存储主机读取从机的数据长度。0：其它模式。（读/写）
\label{fielddesc:SPISLVWRDMABYTELENEN}\item [SPI\_SLV\_WRDMA\_BYTELEN\_EN] 1：DMA 控制模式下 (Wr\_DMA)，SPI\_SLV\_DATA\_BYTELEN 存储主机写入从机的数据长度。0：其它模式。（读/写）
\label{fielddesc:SPISLVRDBUFBYTELENEN}\item [SPI\_SLV\_RDBUF\_BYTELEN\_EN] 1：CPU 控制模式下 (Rd\_BUF)，SPI\_SLV\_DATA\_BYTELEN 存储主机读取从机的数据长度。0：其它模式。（读/写）
\label{fielddesc:SPISLVWRBUFBYTELENEN}\item [SPI\_SLV\_WRBUF\_BYTELEN\_EN] 1：CPU 控制模式下 (Wr\_BUF)，SPI\_SLV\_DATA\_BYTELEN 存储主机写入从机的数据长度。0：其它模式。（读/写）
\label{fielddesc:SPIDMASEGMAGICVALUE}\item [SPI\_DMA\_SEG\_MAGIC\_VALUE] 主机 DMA 分段配置传输模式下，位图表的 Magic 值。（读/写）
\label{fielddesc:SPISLVRDDMADONE}\item [SPI\_SLV\_RD\_DMA\_DONE] 从机模式下，完成 Rd\_DMA 操作将触发中断，该位为中断原始位。CONF\_buf 不可更改该设置。（读/写）
\label{fielddesc:SPIUSRCONF}\item [SPI\_USR\_CONF] 1：使能当前分段配置传输操作的 DMA CONF 阶段，即开始分段配置传输操作。0：使用非分段配置传输模式，即单次传输模式。（读/写）
\end{reglist}\end{regdesc}
\end{register}


\begin{register}{H}{SPI\_DMA\_CONF\_REG}{\hexprefix{}004C}\label{regdesc:SPIDMACONFREG}
\regfield{(reserved)}{3}{29}{000}%
\regfield{SPI\_DMA\_SEG\_TRANS\_CLR}{1}{28}{{0}}%
\regfield{SPI\_EXT\_MEM\_BK\_SIZE}{2}{26}{{0}}%
\regfield{(reserved)}{2}{24}{00}%
\regfield{SPI\_DMA\_OUTFIFO\_EMPTY\_CLR}{1}{23}{{0}}%
\regfield{SPI\_DMA\_INFIFO\_FULL\_CLR}{1}{22}{{0}}%
\regfield{SPI\_RX\_EOF\_EN}{1}{21}{{0}}%
\regfield{SPI\_SLV\_TX\_SEG\_TRANS\_CLR\_EN}{1}{20}{{0}}%
\regfield{SPI\_SLV\_RX\_SEG\_TRANS\_CLR\_EN}{1}{19}{{0}}%
\regfield{SPI\_DMA\_SLV\_SEG\_TRANS\_EN}{1}{18}{{0}}%
\regfield{SPI\_SLV\_LAST\_SEG\_POP\_CLR}{1}{17}{{0}}%
\regfield{SPI\_DMA\_CONTINUE}{1}{16}{{0}}%
\regfield{SPI\_DMA\_TX\_STOP}{1}{15}{{0}}%
\regfield{SPI\_DMA\_RX\_STOP}{1}{14}{{0}}%
\regfield{SPI\_MEM\_TRANS\_EN}{1}{13}{{0}}%
\regfield{SPI\_OUT\_DATA\_BURST\_EN}{1}{12}{{0}}%
\regfield{SPI\_INDSCR\_BURST\_EN}{1}{11}{{0}}%
\regfield{SPI\_OUTDSCR\_BURST\_EN}{1}{10}{{0}}%
\regfield{SPI\_OUT\_EOF\_MODE}{1}{9}{{1}}%
\regfield{(reserved)}{3}{6}{000}%
\regfield{SPI\_AHBM\_RST}{1}{5}{{0}}%
\regfield{SPI\_AHBM\_FIFO\_RST}{1}{4}{{0}}%
\regfield{SPI\_OUT\_RST}{1}{3}{{0}}%
\regfield{SPI\_IN\_RST}{1}{2}{{0}}%
\regfield{(reserved)}{2}{0}{00}%
\reglabel{Reset}\regnewline%
\begin{regdesc}\begin{reglist}
\label{fielddesc:SPIINRST}\item [SPI\_IN\_RST] 该位用于复位 inDMA FSM 和 indata FIFO 指针。（读/写）
\label{fielddesc:SPIOUTRST}\item [SPI\_OUT\_RST] 该位用于复位 outDMA FSM 和 outdata FIFO 指针。（读/写）
\label{fielddesc:SPIAHBMFIFORST}\item [SPI\_AHBM\_FIFO\_RST] 用于复位 SPI DMA AHB 主机 FIFO 指针。（读/写）
\label{fielddesc:SPIAHBMRST}\item [SPI\_AHBM\_RST] 用于复位 SPI DMA AHB 主机。（读/写）
\label{fielddesc:SPIOUTEOFMODE}\item [SPI\_OUT\_EOF\_MODE] OUT\_EOF 标志生成模式。1：DMA 将数据全部从 FIFO 推出后，生成 OUT\_EOF 标志；0：AHB 将数据全部推送到 FIFO 后，生成 OUT\_EOF 标志。（读/写）
\label{fielddesc:SPIOUTDSCRBURSTEN}\item [SPI\_OUTDSCR\_BURST\_EN] 从内存读取数据时，读取描述符将使用 burst 模式。（读/写）
\label{fielddesc:SPIINDSCRBURSTEN}\item [SPI\_INDSCR\_BURST\_EN] 数据写入内存时，读取描述符将使用 burst 模式。（读/写）
\label{fielddesc:SPIOUTDATABURSTEN}\item [SPI\_OUT\_DATA\_BURST\_EN] SPI DMA 使用 burst 模式从内存读取数据。（读/写）
\label{fielddesc:SPIMEMTRANSEN}\item [SPI\_MEM\_TRANS\_EN] 1：内存数据传输使能位。将 DMA RX buffer 数据发送到 DMA TX buffer。0：禁用此功能。（读/写）
\label{fielddesc:SPIDMARXSTOP}\item [SPI\_DMA\_RX\_STOP] SPI DMA 连续发送或接收模式下，SPI DMA 读数据停止。（读/写）
\label{fielddesc:SPIDMATXSTOP}\item [SPI\_DMA\_TX\_STOP] SPI DMA 继续发送或接收模式下，SPI DMA 写数据停止。（读/写）
\label{fielddesc:SPIDMACONTINUE}\item [SPI\_DMA\_CONTINUE] SPI DMA 继续发送或接收数据。（读/写）
\label{fielddesc:SPISLVLASTSEGPOPCLR}\item [SPI\_SLV\_LAST\_SEG\_POP\_CLR] 在从机连续传输前置位此位，然后再清零此位，就可以用于清除上次传输中未使用的 DMA TX 数据。（读/写）
\label{fielddesc:SPIDMASLVSEGTRANSEN}\item [SPI\_DMA\_SLV\_SEG\_TRANS\_EN] SPI DMA 半双工从机模式下，使能连续传输功能。1：使能此功能；0：禁用此功能。（读/写）
\label{fielddesc:SPISLVRXSEGTRANSCLREN}\item [SPI\_SLV\_RX\_SEG\_TRANS\_CLR\_EN] SPI DMA 半双工从机模式下，当 DMA RX buffer 小于实际接收的数据长度时，1：后续传输的数据都不接收；0：本次传输的数据不接收，下次传输时, 如果 DMA RX Buffer 长度不为零，则继续接收，否则不接收。（读/写）
\label{fielddesc:SPISLVTXSEGTRANSCLREN}\item [SPI\_SLV\_TX\_SEG\_TRANS\_CLR\_EN] SPI DMA 半双工从机模式下，当 DMA TX buffer 大小小于实际发送的数据长度时，1：后续传输的数据都不更新，发送同一个旧数据；0：本次传输的数据都不更新，下次传输时, 如果 DMA TX FIFO 填充了新的数据，则继续发送新数据，否则发送数据不更新。（读/写）
\item [见下页]
\end{reglist}\end{regdesc}
\end{register}


\addtocounter{Regfloat}{-1}
\begin{register}{H}{SPI\_DMA\_CONF\_REG}{\hexprefix{}004C}
\begin{regdesc}\begin{reglist}
\item [接上页]
\label{fielddesc:SPIRXEOFEN}\item [SPI\_RX\_EOF\_EN] 1：在 SPI DMA 数据传输过程中，如果 DMA 传输的数据字节数等于 SPI\_SLV\_DMA\_RD\_BYTELEN[19:0] 或 SPI\_MST\_DMA\_RD\_BYTELEN[19:0]，则置位 SPI\_IN\_SUC\_EOF\_INT\_RAW。0：在非分段配置传输模式下，SPI\_IN\_SUC\_EOF\_INT\_RAW 由 SPI\_TRANS\_DONE 置位；或在分段配置传输模式下，由 SPI\_DMA\_SEG\_TRANS\_DONE 置位。（读/写）
\label{fielddesc:SPIDMAINFIFOFULLCLR}\item [SPI\_DMA\_INFIFO\_FULL\_CLR] 当置位 SPI\_SLV\_RX\_SEG\_TRANS\_CLR\_EN 比特，且 SPI DMA 半双工从机模式下，出现 DMA RX buffer 大小小于传输数据的长度时，需置位  [SPI\_DMA\_INFIFO\_FULL\_CLR]  比特，并清零该比特，从而不影响下次正常传输。（读/写）
\label{fielddesc:SPIDMAOUTFIFOEMPTYCLR}\item [SPI\_DMA\_OUTFIFO\_EMPTY\_CLR]
当置位 SPI\_SLV\_TX\_SEG\_TRANS\_CLR\_EN 比特，且 SPI DMA 半双工从机模式下，出现 DMA TX buffer 大小小于传输数据的长度，需置位  [SPI\_DMA\_OUTFIFO\_EMPTY\_CLR]  比特，并清零该比特，从而不影响下次正常传输。（读/写）
\label{fielddesc:SPIEXTMEMBKSIZE}\item [SPI\_EXT\_MEM\_BK\_SIZE] 选择外部存储器的大小。（读/写）
\label{fielddesc:SPIDMASEGTRANSCLR}\item [SPI\_DMA\_SEG\_TRANS\_CLR] 1：结束从机连续传输模式，相当于 \hexprefix{}05 命令。两个及两个以上的连续传输结束信号可能会在 DMA RX 引入错误。0：使用其它模式。一个 APB CLK 周期后将由硬件清零。（读/写）
\end{reglist}\end{regdesc}
\end{register}


\begin{register}{H}{SPI\_DMA\_OUT\_LINK\_REG}{\hexprefix{}0050}\label{regdesc:SPIDMAOUTLINKREG}
\regfield{SPI\_DMA\_TX\_ENA}{1}{31}{{0}}%
\regfield{SPI\_OUTLINK\_RESTART}{1}{30}{{0}}%
\regfield{SPI\_OUTLINK\_START}{1}{29}{{0}}%
\regfield{SPI\_OUTLINK\_STOP}{1}{28}{{0}}%
\regfield{(reserved)}{8}{20}{00000000}%
\regfield{SPI\_OUTLINK\_ADDR}{20}{0}{{\hexprefix{}000}}%
\reglabel{Reset}\regnewline%
\begin{regdesc}\begin{reglist}
\label{fielddesc:SPIOUTLINKADDR}\item [SPI\_OUTLINK\_ADDR] 第一个 outlink 描述符地址。（读/写）
\label{fielddesc:SPIOUTLINKSTOP}\item [SPI\_OUTLINK\_STOP] 置位此位，停止使用 outlink 描述符。（读/写）
\label{fielddesc:SPIOUTLINKSTART}\item [SPI\_OUTLINK\_START] 置位此位，开始使用 outlink 描述符。（读/写）
\label{fielddesc:SPIOUTLINKRESTART}\item [SPI\_OUTLINK\_RESTART] 置位此位，挂载新的 outlink 描述符。（读/写）
\label{fielddesc:SPIDMATXENA}\item [SPI\_DMA\_TX\_ENA] 置位此位，使能 DMA 控制的发送数据模式；清零此位，使能 CPU 控制的发送数据模式。（读/写）
\end{reglist}\end{regdesc}
\end{register}


\begin{register}{H}{SPI\_IN\_ERR\_EOF\_DES\_ADDR\_REG}{\hexprefix{}0068}\label{regdesc:SPIINERREOFDESADDRREG}
\regfield{SPI\_DMA\_IN\_ERR\_EOF\_DES\_ADDR}{32}{0}{{0}}%
\reglabel{Reset}\regnewline%
\begin{regdesc}\begin{reglist}
\label{fielddesc:SPIDMAINERREOFDESADDR}\item [SPI\_DMA\_IN\_ERR\_EOF\_DES\_ADDR] SPI DMA 产生接收错误时，inlink 描述符的地址。 （只读）
\end{reglist}\end{regdesc}
\end{register}


\begin{register}{H}{SPI\_IN\_SUC\_EOF\_DES\_ADDR\_REG}{\hexprefix{}006C}\label{regdesc:SPIINSUCEOFDESADDRREG}
\regfield{SPI\_DMA\_IN\_SUC\_EOF\_DES\_ADDR}{32}{0}{{0}}%
\reglabel{Reset}\regnewline%
\begin{regdesc}\begin{reglist}
\label{fielddesc:SPIDMAINSUCEOFDESADDR}\item [SPI\_DMA\_IN\_SUC\_EOF\_DES\_ADDR] SPI DMA 产生 FROM\_SUC\_EOF 时，上一次 inlink 描述符的地址。 （只读）
\end{reglist}\end{regdesc}
\end{register}


\begin{register}{H}{SPI\_INLINK\_DSCR\_REG}{\hexprefix{}0070}\label{regdesc:SPIINLINKDSCRREG}
\regfield{SPI\_DMA\_INLINK\_DSCR}{32}{0}{{0}}%
\reglabel{Reset}\regnewline%
\begin{regdesc}\begin{reglist}
\label{fielddesc:SPIDMAINLINKDSCR}\item [SPI\_DMA\_INLINK\_DSCR] 当前 inlink 描述符指针的内容。 （只读）
\end{reglist}\end{regdesc}
\end{register}


\begin{register}{H}{SPI\_INLINK\_DSCR\_BF0\_REG}{\hexprefix{}0074}\label{regdesc:SPIINLINKDSCRBF0REG}
\regfield{SPI\_DMA\_INLINK\_DSCR\_BF0}{32}{0}{{0}}%
\reglabel{Reset}\regnewline%
\begin{regdesc}\begin{reglist}
\label{fielddesc:SPIDMAINLINKDSCRBF0}\item [SPI\_DMA\_INLINK\_DSCR\_BF0] 下个 inlink 描述符指针的内容。 （只读）
\end{reglist}\end{regdesc}
\end{register}


\begin{register}{H}{SPI\_OUT\_EOF\_BFR\_DES\_ADDR\_REG}{\hexprefix{}007C}\label{regdesc:SPIOUTEOFBFRDESADDRREG}
\regfield{SPI\_DMA\_OUT\_EOF\_BFR\_DES\_ADDR}{32}{0}{{0}}%
\reglabel{Reset}\regnewline%
\begin{regdesc}\begin{reglist}
\label{fielddesc:SPIDMAOUTEOFBFRDESADDR}\item [SPI\_DMA\_OUT\_EOF\_BFR\_DES\_ADDR] 生成 EOF 的发送链表描述符对应的 buffer 地址。（只读）
\end{reglist}\end{regdesc}
\end{register}


\begin{register}{H}{SPI\_OUT\_EOF\_DES\_ADDR\_REG}{\hexprefix{}0080}\label{regdesc:SPIOUTEOFDESADDRREG}
\regfield{SPI\_DMA\_OUT\_EOF\_DES\_ADDR}{32}{0}{{0}}%
\reglabel{Reset}\regnewline%
\begin{regdesc}\begin{reglist}
\label{fielddesc:SPIDMAOUTEOFDESADDR}\item [SPI\_DMA\_OUT\_EOF\_DES\_ADDR] SPI DMA 生成 TO\_EOF 时，上次 outlink 描述符的地址。（只读）
\end{reglist}\end{regdesc}
\end{register}


\begin{register}{H}{SPI\_OUTLINK\_DSCR\_REG}{\hexprefix{}0084}\label{regdesc:SPIOUTLINKDSCRREG}
\regfield{SPI\_DMA\_OUTLINK\_DSCR}{32}{0}{{0}}%
\reglabel{Reset}\regnewline%
\begin{regdesc}\begin{reglist}
\label{fielddesc:SPIDMAOUTLINKDSCR}\item [SPI\_DMA\_OUTLINK\_DSCR] 当前 outlink 描述符指针的内容。 （只读）
\end{reglist}\end{regdesc}
\end{register}


\begin{register}{H}{SPI\_OUTLINK\_DSCR\_BF0\_REG}{\hexprefix{}0088}\label{regdesc:SPIOUTLINKDSCRBF0REG}
\regfield{SPI\_DMA\_OUTLINK\_DSCR\_BF0}{32}{0}{{0}}%
\reglabel{Reset}\regnewline%
\begin{regdesc}\begin{reglist}
\label{fielddesc:SPIDMAOUTLINKDSCRBF0}\item [SPI\_DMA\_OUTLINK\_DSCR\_BF0] 下个 outlink 描述符指针的内容。 （只读）
\end{reglist}\end{regdesc}
\end{register}


\begin{register}{H}{SPI\_DMA\_OUTSTATUS\_REG}{\hexprefix{}0090}\label{regdesc:SPIDMAOUTSTATUSREG}
\regfield{SPI\_DMA\_OUTFIFO\_EMPTY}{1}{31}{{1}}%
\regfield{SPI\_DMA\_OUTFIFO\_FULL}{1}{30}{{0}}%
\regfield{SPI\_DMA\_OUTFIFO\_CNT}{7}{23}{{0}}%
\regfield{SPI\_DMA\_OUT\_STATE}{3}{20}{{0}}%
\regfield{SPI\_DMA\_OUTDSCR\_STATE}{2}{18}{{0}}%
\regfield{SPI\_DMA\_OUTDSCR\_ADDR}{18}{0}{{0}}%
\reglabel{Reset}\regnewline%
\begin{regdesc}\begin{reglist}
\label{fielddesc:SPIDMAOUTDSCRADDR}\item [SPI\_DMA\_OUTDSCR\_ADDR] SPI DMA out 描述符地址。 （只读）
\label{fielddesc:SPIDMAOUTDSCRSTATE}\item [SPI\_DMA\_OUTDSCR\_STATE] SPI DMA out 描述符状态。 （只读）
\label{fielddesc:SPIDMAOUTSTATE}\item [SPI\_DMA\_OUT\_STATE] SPI DMA out 数据状态。 （只读）
\label{fielddesc:SPIDMAOUTFIFOCNT}\item [SPI\_DMA\_OUTFIFO\_CNT] SPI DMA outFIFO 剩余数据。 （只读）
\label{fielddesc:SPIDMAOUTFIFOFULL}\item [SPI\_DMA\_OUTFIFO\_FULL] SPI DMA outFIFO 已满。 （只读）
\label{fielddesc:SPIDMAOUTFIFOEMPTY}\item [SPI\_DMA\_OUTFIFO\_EMPTY] SPI DMA outFIFO 已空。 （只读）
\end{reglist}\end{regdesc}
\end{register}


\begin{register}{H}{SPI\_DMA\_INSTATUS\_REG}{\hexprefix{}0094}\label{regdesc:SPIDMAINSTATUSREG}
\regfield{SPI\_DMA\_INFIFO\_EMPTY}{1}{31}{{1}}%
\regfield{SPI\_DMA\_INFIFO\_FULL}{1}{30}{{0}}%
\regfield{SPI\_DMA\_INFIFO\_CNT}{7}{23}{{0}}%
\regfield{SPI\_DMA\_IN\_STATE}{3}{20}{{0}}%
\regfield{SPI\_DMA\_INDSCR\_STATE}{2}{18}{{0}}%
\regfield{SPI\_DMA\_INDSCR\_ADDR}{18}{0}{{0}}%
\reglabel{Reset}\regnewline%
\begin{regdesc}\begin{reglist}
\label{fielddesc:SPIDMAINDSCRADDR}\item [SPI\_DMA\_INDSCR\_ADDR] SPI DMA in 描述符地址。 （只读）
\label{fielddesc:SPIDMAINDSCRSTATE}\item [SPI\_DMA\_INDSCR\_STATE] SPI DMA in 描述符状态。 （只读）
\label{fielddesc:SPIDMAINSTATE}\item [SPI\_DMA\_IN\_STATE] SPI DMA in 数据状态。 （只读）
\label{fielddesc:SPIDMAINFIFOCNT}\item [SPI\_DMA\_INFIFO\_CNT] SPI DMA inFIFO 剩余数据。 （只读）
\label{fielddesc:SPIDMAINFIFOFULL}\item [SPI\_DMA\_INFIFO\_FULL] SPI DMA inFIFO 已满。 （只读）
\label{fielddesc:SPIDMAINFIFOEMPTY}\item [SPI\_DMA\_INFIFO\_EMPTY] SPI DMA inFIFO 已空。 （只读）
\end{reglist}\end{regdesc}
\end{register}


\begin{register}{H}{SPI\_DMA\_IN\_LINK\_REG}{\hexprefix{}0054}\label{regdesc:SPIDMAINLINKREG}
\regfield{SPI\_DMA\_RX\_ENA}{1}{31}{{0}}%
\regfield{SPI\_INLINK\_RESTART}{1}{30}{{0}}%
\regfield{SPI\_INLINK\_START}{1}{29}{{0}}%
\regfield{SPI\_INLINK\_STOP}{1}{28}{{0}}%
\regfield{(reserved)}{7}{21}{0000000}%
\regfield{SPI\_INLINK\_AUTO\_RET}{1}{20}{{0}}%
\regfield{SPI\_INLINK\_ADDR}{20}{0}{{\hexprefix{}000}}%
\reglabel{Reset}\regnewline%
\begin{regdesc}\begin{reglist}
\label{fielddesc:SPIINLINKADDR}\item [SPI\_INLINK\_ADDR] 第一个 inlink 描述符地址。（读/写）
\label{fielddesc:SPIINLINKAUTORET}\item [SPI\_INLINK\_AUTO\_RET] 置位此位，则当数据包出错时，inlink 描述符将返回到第一个 link 节点。（读/写）
\label{fielddesc:SPIINLINKSTOP}\item [SPI\_INLINK\_STOP] 置位此位，停止使用 inlink 描述符。（读/写）
\label{fielddesc:SPIINLINKSTART}\item [SPI\_INLINK\_START] 置位此位，开始使用 inlink 描述符。（读/写）
\label{fielddesc:SPIINLINKRESTART}\item [SPI\_INLINK\_RESTART] 置位此位，挂载新的 inlink 描述符。（读/写）
\label{fielddesc:SPIDMARXENA}\item [SPI\_DMA\_RX\_ENA] 置位此位，使能 DMA 控制的接收数据模式；清零此位，使能 CPU 控制的接收数据模式。（读/写）
\end{reglist}\end{regdesc}
\end{register}


\begin{register}{H}{SPI\_DMA\_INT\_ENA\_REG}{\hexprefix{}0058}\label{regdesc:SPIDMAINTENAREG}
\regfield{(reserved)}{16}{16}{0000000000000000}%
\regfield{SPI\_SLV\_CMDA\_INT\_ENA}{1}{15}{{0}}%
\regfield{SPI\_SLV\_CMD9\_INT\_ENA}{1}{14}{{0}}%
\regfield{SPI\_SLV\_CMD8\_INT\_ENA}{1}{13}{{0}}%
\regfield{SPI\_SLV\_CMD7\_INT\_ENA}{1}{12}{{0}}%
\regfield{SPI\_SLV\_CMD6\_INT\_ENA}{1}{11}{{0}}%
\regfield{SPI\_OUTFIFO\_EMPTY\_ERR\_INT\_ENA}{1}{10}{{0}}%
\regfield{SPI\_INFIFO\_FULL\_ERR\_INT\_ENA}{1}{9}{{0}}%
\regfield{SPI\_OUT\_TOTAL\_EOF\_INT\_ENA}{1}{8}{{0}}%
\regfield{SPI\_OUT\_EOF\_INT\_ENA}{1}{7}{{0}}%
\regfield{SPI\_OUT\_DONE\_INT\_ENA}{1}{6}{{0}}%
\regfield{SPI\_IN\_SUC\_EOF\_INT\_ENA}{1}{5}{{0}}%
\regfield{SPI\_IN\_ERR\_EOF\_INT\_ENA}{1}{4}{{0}}%
\regfield{SPI\_IN\_DONE\_INT\_ENA}{1}{3}{{0}}%
\regfield{SPI\_INLINK\_DSCR\_ERROR\_INT\_ENA}{1}{2}{{0}}%
\regfield{SPI\_OUTLINK\_DSCR\_ERROR\_INT\_ENA}{1}{1}{{0}}%
\regfield{SPI\_INLINK\_DSCR\_EMPTY\_INT\_ENA}{1}{0}{{0}}%
\reglabel{Reset}\regnewline%
\begin{regdesc}\begin{reglist}
\label{fielddesc:SPIINLINKDSCREMPTYINTENA}\item [SPI\_INLINK\_DSCR\_EMPTY\_INT\_ENA] Inlink 描述符不足将触发中断，此位为该中断的使能位。可在 CONF 阶段配置。（读/写）
\label{fielddesc:SPIOUTLINKDSCRERRORINTENA}\item [SPI\_OUTLINK\_DSCR\_ERROR\_INT\_ENA] Outlink 描述符出现错误将触发中断，此位为该中断的使能位。可在 CONF 阶段配置。（读/写）
\label{fielddesc:SPIINLINKDSCRERRORINTENA}\item [SPI\_INLINK\_DSCR\_ERROR\_INT\_ENA] Inlink 描述符出现错误将触发中断，此位为该中断的使能位。可在 CONF 阶段配置。（读/写）
\label{fielddesc:SPIINDONEINTENA}\item [SPI\_IN\_DONE\_INT\_ENA] 一个 inlink 描述符使用完成后将触发中断，此位为该中断的使能位。可在 CONF 阶段配置。（读/写）
\label{fielddesc:SPIINERREOFINTENA}\item [SPI\_IN\_ERR\_EOF\_INT\_ENA] 接收出现错误将触发中断，此位为该中断的使能位。可在 CONF 阶段配置。（读/写）
\label{fielddesc:SPIINSUCEOFINTENA}\item [SPI\_IN\_SUC\_EOF\_INT\_ENA] 所有数据包全部从主机接收后将触发中断，此位为该中断的使能位。可在 CONF 阶段配置。（读/写）
\label{fielddesc:SPIOUTDONEINTENA}\item [SPI\_OUT\_DONE\_INT\_ENA] 一个 outlink 描述符使用完成后将触发中断，此位为该中断的使能位。可在 CONF 阶段配置。（读/写）
\label{fielddesc:SPIOUTEOFINTENA}\item [SPI\_OUT\_EOF\_INT\_ENA] 单个数据包发送到主机后将触发中断，此位为该中断的使能位。可在 CONF 阶段配置。（读/写）
\label{fielddesc:SPIOUTTOTALEOFINTENA}\item [SPI\_OUT\_TOTAL\_EOF\_INT\_ENA] 所有数据包全部发送到主机后将触发中断，此位为该中断的使能位。可在 CONF 阶段配置。（读/写）
\label{fielddesc:SPIINFIFOFULLERRINTENA}\item [SPI\_INFIFO\_FULL\_ERR\_INT\_ENA] Infifo 满错误中断的使能位。（读/写）
\label{fielddesc:SPIOUTFIFOEMPTYERRINTENA}\item [SPI\_OUTFIFO\_EMPTY\_ERR\_INT\_ENA] Outfifo 空错误中断的使能位。（读/写）
\label{fielddesc:SPISLVCMD6INTENA}\item [SPI\_SLV\_CMD6\_INT\_ENA] SPI 从机 CMD6 中断使能位。（读/写）
\label{fielddesc:SPISLVCMD7INTENA}\item [SPI\_SLV\_CMD7\_INT\_ENA] SPI 从机 CMD7 中断使能位。（读/写）
\label{fielddesc:SPISLVCMD8INTENA}\item [SPI\_SLV\_CMD8\_INT\_ENA] SPI 从机 CMD8 中断使能位。（读/写）
\label{fielddesc:SPISLVCMD9INTENA}\item [SPI\_SLV\_CMD9\_INT\_ENA] SPI 从机 CMD9 中断使能位。（读/写）
\label{fielddesc:SPISLVCMDAINTENA}\item [SPI\_SLV\_CMDA\_INT\_ENA] SPI 从机 CMDA 中断使能位。（读/写）
\end{reglist}\end{regdesc}
\end{register}


\begin{register}{H}{SPI\_DMA\_INT\_RAW\_REG}{\hexprefix{}005C}\label{regdesc:SPIDMAINTRAWREG}
\regfield{(reserved)}{16}{16}{0000000000000000}%
\regfield{SPI\_SLV\_CMDA\_INT\_RAW}{1}{15}{{0}}%
\regfield{SPI\_SLV\_CMD9\_INT\_RAW}{1}{14}{{0}}%
\regfield{SPI\_SLV\_CMD8\_INT\_RAW}{1}{13}{{0}}%
\regfield{SPI\_SLV\_CMD7\_INT\_RAW}{1}{12}{{0}}%
\regfield{SPI\_SLV\_CMD6\_INT\_RAW}{1}{11}{{0}}%
\regfield{SPI\_OUTFIFO\_EMPTY\_ERR\_INT\_RAW}{1}{10}{{0}}%
\regfield{SPI\_INFIFO\_FULL\_ERR\_INT\_RAW}{1}{9}{{0}}%
\regfield{SPI\_OUT\_TOTAL\_EOF\_INT\_RAW}{1}{8}{{0}}%
\regfield{SPI\_OUT\_EOF\_INT\_RAW}{1}{7}{{0}}%
\regfield{SPI\_OUT\_DONE\_INT\_RAW}{1}{6}{{0}}%
\regfield{SPI\_IN\_SUC\_EOF\_INT\_RAW}{1}{5}{{0}}%
\regfield{SPI\_IN\_ERR\_EOF\_INT\_RAW}{1}{4}{{0}}%
\regfield{SPI\_IN\_DONE\_INT\_RAW}{1}{3}{{0}}%
\regfield{SPI\_INLINK\_DSCR\_ERROR\_INT\_RAW}{1}{2}{{0}}%
\regfield{SPI\_OUTLINK\_DSCR\_ERROR\_INT\_RAW}{1}{1}{{0}}%
\regfield{SPI\_INLINK\_DSCR\_EMPTY\_INT\_RAW}{1}{0}{{0}}%
\reglabel{Reset}\regnewline%
\begin{regdesc}\begin{reglist}
\label{fielddesc:SPIINLINKDSCREMPTYINTRAW}\item [SPI\_INLINK\_DSCR\_EMPTY\_INT\_RAW] Inlink 描述符不足将触发中断，此位为该中断的原始位。可在 CONF 阶段配置。 （只读）
\label{fielddesc:SPIOUTLINKDSCRERRORINTRAW}\item [SPI\_OUTLINK\_DSCR\_ERROR\_INT\_RAW] Outlink 描述符出现错误将触发中断，此位为该中断的原始位。可在 CONF 阶段配置。 （只读）
\label{fielddesc:SPIINLINKDSCRERRORINTRAW}\item [SPI\_INLINK\_DSCR\_ERROR\_INT\_RAW] Inlink 描述符出现错误将触发中断，此位为该中断的原始位。可在 CONF 阶段配置。 （只读）
\label{fielddesc:SPIINDONEINTRAW}\item [SPI\_IN\_DONE\_INT\_RAW] 一个 inlink 描述符使用完成后将触发中断，此位为该中断的原始位。可在 CONF 阶段配置。 （只读）
\label{fielddesc:SPIINERREOFINTRAW}\item [SPI\_IN\_ERR\_EOF\_INT\_RAW] 接收出现错误将触发中断，此位为该中断的原始位。可在 CONF 阶段配置。 （只读）
\label{fielddesc:SPIINSUCEOFINTRAW}\item [SPI\_IN\_SUC\_EOF\_INT\_RAW] 所有数据包全部从主机接收后将触发中断，此位为该中断的原始位。可在 CONF 阶段配置。 （只读）
\label{fielddesc:SPIOUTDONEINTRAW}\item [SPI\_OUT\_DONE\_INT\_RAW] 一个 outlink 描述符使用完成后将触发中断，此位为该中断的原始位。可在 CONF 阶段配置。 （只读）
\label{fielddesc:SPIOUTEOFINTRAW}\item [SPI\_OUT\_EOF\_INT\_RAW] 单个数据包发送到主机后将触发中断，此位为该中断的原始位。可在 CONF 阶段配置。 （只读）
\label{fielddesc:SPIOUTTOTALEOFINTRAW}\item [SPI\_OUT\_TOTAL\_EOF\_INT\_RAW] 所有数据包全部发送到主机后将触发中断，此位为该中断的原始位。可在 CONF 阶段配置。 （只读）
\label{fielddesc:SPIINFIFOFULLERRINTRAW}\item [SPI\_INFIFO\_FULL\_ERR\_INT\_RAW] 当 DMA RX buffer 的大小小于传输数据长度时，触发此中断。 不可在 CONF 阶段配置。 （只读）
\label{fielddesc:SPIOUTFIFOEMPTYERRINTRAW}\item [SPI\_OUTFIFO\_EMPTY\_ERR\_INT\_RAW] 当 DMA TX buffer 的大小小于传输数据长度时，触发此中断。不可在 CONF 阶段配置。（只读）
\label{fielddesc:SPISLVCMD6INTRAW}\item [SPI\_SLV\_CMD6\_INT\_RAW] SPI 从机 CMD6 中断原始位。（读/写）
\label{fielddesc:SPISLVCMD7INTRAW}\item [SPI\_SLV\_CMD7\_INT\_RAW] SPI 从机 CMD7 中断原始位。（读/写）
\label{fielddesc:SPISLVCMD8INTRAW}\item [SPI\_SLV\_CMD8\_INT\_RAW] SPI 从机 CMD8 中断原始位。（读/写）
\label{fielddesc:SPISLVCMD9INTRAW}\item [SPI\_SLV\_CMD9\_INT\_RAW] SPI 从机 CMD9 中断原始位。（读/写）
\label{fielddesc:SPISLVCMDAINTRAW}\item [SPI\_SLV\_CMDA\_INT\_RAW] SPI 从机 CMDA 中断原始位。（读/写）
\end{reglist}\end{regdesc}
\end{register}


\begin{register}{H}{SPI\_DMA\_INT\_ST\_REG}{\hexprefix{}0060}\label{regdesc:SPIDMAINTSTREG}
\regfield{(reserved)}{16}{16}{0000000000000000}%
\regfield{SPI\_SLV\_CMDA\_INT\_ST}{1}{15}{{0}}%
\regfield{SPI\_SLV\_CMD9\_INT\_ST}{1}{14}{{0}}%
\regfield{SPI\_SLV\_CMD8\_INT\_ST}{1}{13}{{0}}%
\regfield{SPI\_SLV\_CMD7\_INT\_ST}{1}{12}{{0}}%
\regfield{SPI\_SLV\_CMD6\_INT\_ST}{1}{11}{{0}}%
\regfield{SPI\_OUTFIFO\_EMPTY\_ERR\_INT\_ST}{1}{10}{{0}}%
\regfield{SPI\_INFIFO\_FULL\_ERR\_INT\_ST}{1}{9}{{0}}%
\regfield{SPI\_OUT\_TOTAL\_EOF\_INT\_ST}{1}{8}{{0}}%
\regfield{SPI\_OUT\_EOF\_INT\_ST}{1}{7}{{0}}%
\regfield{SPI\_OUT\_DONE\_INT\_ST}{1}{6}{{0}}%
\regfield{SPI\_IN\_SUC\_EOF\_INT\_ST}{1}{5}{{0}}%
\regfield{SPI\_IN\_ERR\_EOF\_INT\_ST}{1}{4}{{0}}%
\regfield{SPI\_IN\_DONE\_INT\_ST}{1}{3}{{0}}%
\regfield{SPI\_INLINK\_DSCR\_ERROR\_INT\_ST}{1}{2}{{0}}%
\regfield{SPI\_OUTLINK\_DSCR\_ERROR\_INT\_ST}{1}{1}{{0}}%
\regfield{SPI\_INLINK\_DSCR\_EMPTY\_INT\_ST}{1}{0}{{0}}%
\reglabel{Reset}\regnewline%
\begin{regdesc}\begin{reglist}
\label{fielddesc:SPIINLINKDSCREMPTYINTST}\item [SPI\_INLINK\_DSCR\_EMPTY\_INT\_ST] Inlink 描述符不足将触发中断，此位为该中断的状态位。 （只读）
\label{fielddesc:SPIOUTLINKDSCRERRORINTST}\item [SPI\_OUTLINK\_DSCR\_ERROR\_INT\_ST] Outlink 描述符错误中断的状态位。 （只读）
\label{fielddesc:SPIINLINKDSCRERRORINTST}\item [SPI\_INLINK\_DSCR\_ERROR\_INT\_ST] Inlink 描述符错误中断的状态位。 （只读）
\label{fielddesc:SPIINDONEINTST}\item [SPI\_IN\_DONE\_INT\_ST] 一个 inlink 描述符使用完成后将触发中断，此位为该中断的状态位。 （只读）
\label{fielddesc:SPIINERREOFINTST}\item [SPI\_IN\_ERR\_EOF\_INT\_ST] 接收错误中断的状态位。 （只读）
\label{fielddesc:SPIINSUCEOFINTST}\item [SPI\_IN\_SUC\_EOF\_INT\_ST] 所有数据包全部从主机接收后将触发中断，此位为该中断的状态位。 （只读）
\label{fielddesc:SPIOUTDONEINTST}\item [SPI\_OUT\_DONE\_INT\_ST] 一个 outlink 描述符使用完成后将触发中断，此位为该中断的状态位。 （只读）
\label{fielddesc:SPIOUTEOFINTST}\item [SPI\_OUT\_EOF\_INT\_ST] 单个数据包发送到主机后将触发中断，此位为该中断的状态位。 （只读）
\label{fielddesc:SPIOUTTOTALEOFINTST}\item [SPI\_OUT\_TOTAL\_EOF\_INT\_ST] 所有数据包全部发送到主机后将触发中断，此位为该中断的状态位。 （只读）
\label{fielddesc:SPIINFIFOFULLERRINTST}\item [SPI\_INFIFO\_FULL\_ERR\_INT\_ST] Infifo 完整错误中断的状态位。 （只读）
\label{fielddesc:SPIOUTFIFOEMPTYERRINTST}\item [SPI\_OUTFIFO\_EMPTY\_ERR\_INT\_ST] Outfifo 空错误中断的状态位。 （只读）
\label{fielddesc:SPISLVCMD6INTST}\item [SPI\_SLV\_CMD6\_INT\_ST] SPI 从机 CMD6 中断状态位。（读/写）
\label{fielddesc:SPISLVCMD7INTST}\item [SPI\_SLV\_CMD7\_INT\_ST] SPI 从机 CMD7 中断状态位。（读/写）
\label{fielddesc:SPISLVCMD8INTST}\item [SPI\_SLV\_CMD8\_INT\_ST] SPI 从机 CMD8 中断状态位。（读/写）
\label{fielddesc:SPISLVCMD9INTST}\item [SPI\_SLV\_CMD9\_INT\_ST] SPI 从机 CMD9 中断状态位。（读/写）
\label{fielddesc:SPISLVCMDAINTST}\item [SPI\_SLV\_CMDA\_INT\_ST] SPI 从机 CMDA 中断状态位。（读/写）
\end{reglist}\end{regdesc}
\end{register}


\begin{register}{H}{SPI\_DMA\_INT\_CLR\_REG}{\hexprefix{}0064}\label{regdesc:SPIDMAINTCLRREG}
\regfield{(reserved)}{16}{16}{0000000000000000}%
\regfield{SPI\_SLV\_CMDA\_INT\_CLR}{1}{15}{{0}}%
\regfield{SPI\_SLV\_CMD9\_INT\_CLR}{1}{14}{{0}}%
\regfield{SPI\_SLV\_CMD8\_INT\_CLR}{1}{13}{{0}}%
\regfield{SPI\_SLV\_CMD7\_INT\_CLR}{1}{12}{{0}}%
\regfield{SPI\_SLV\_CMD6\_INT\_CLR}{1}{11}{{0}}%
\regfield{SPI\_OUTFIFO\_EMPTY\_ERR\_INT\_CLR}{1}{10}{{0}}%
\regfield{SPI\_INFIFO\_FULL\_ERR\_INT\_CLR}{1}{9}{{0}}%
\regfield{SPI\_OUT\_TOTAL\_EOF\_INT\_CLR}{1}{8}{{0}}%
\regfield{SPI\_OUT\_EOF\_INT\_CLR}{1}{7}{{0}}%
\regfield{SPI\_OUT\_DONE\_INT\_CLR}{1}{6}{{0}}%
\regfield{SPI\_IN\_SUC\_EOF\_INT\_CLR}{1}{5}{{0}}%
\regfield{SPI\_IN\_ERR\_EOF\_INT\_CLR}{1}{4}{{0}}%
\regfield{SPI\_IN\_DONE\_INT\_CLR}{1}{3}{{0}}%
\regfield{SPI\_INLINK\_DSCR\_ERROR\_INT\_CLR}{1}{2}{{0}}%
\regfield{SPI\_OUTLINK\_DSCR\_ERROR\_INT\_CLR}{1}{1}{{0}}%
\regfield{SPI\_INLINK\_DSCR\_EMPTY\_INT\_CLR}{1}{0}{{0}}%
\reglabel{Reset}\regnewline%
\begin{regdesc}\begin{reglist}
\label{fielddesc:SPIINLINKDSCREMPTYINTCLR}\item [SPI\_INLINK\_DSCR\_EMPTY\_INT\_CLR] Inlink 描述符不足将触发中断，此位为该中断的清除位。可在 CONF 状态下配置。（读/写）
\label{fielddesc:SPIOUTLINKDSCRERRORINTCLR}\item [SPI\_OUTLINK\_DSCR\_ERROR\_INT\_CLR] Outlink 描述符出现错误将触发中断，此位为该中断的清除位。可在 CONF 状态下配置。（读/写）
\label{fielddesc:SPIINLINKDSCRERRORINTCLR}\item [SPI\_INLINK\_DSCR\_ERROR\_INT\_CLR] Inlink 描述符出现错误将触发中断，此位为该中断的清除位。可在 CONF 状态下配置。（读/写）
\label{fielddesc:SPIINDONEINTCLR}\item [SPI\_IN\_DONE\_INT\_CLR] 一个 inlink 描述符使用完成后将触发中断，此位为该中断的清除位。可在 CONF 状态下配置。（读/写）
\label{fielddesc:SPIINERREOFINTCLR}\item [SPI\_IN\_ERR\_EOF\_INT\_CLR] 接收出现错误将触发中断，此位为该中断的清除位。可在 CONF 状态下配置。（读/写）
\label{fielddesc:SPIINSUCEOFINTCLR}\item [SPI\_IN\_SUC\_EOF\_INT\_CLR] 所有数据包全部从主机接收后将触发中断，此位为该中断的清除位。可在 CONF 状态下配置。（读/写）
\label{fielddesc:SPIOUTDONEINTCLR}\item [SPI\_OUT\_DONE\_INT\_CLR] 一个 outlink 描述符使用完成后将触发中断，此位为该中断的清除位。可在 CONF 状态下配置。（读/写）
\label{fielddesc:SPIOUTEOFINTCLR}\item [SPI\_OUT\_EOF\_INT\_CLR] 单个数据包发送到主机后将触发中断，此位为该中断的清除位。可在 CONF 状态下配置。（读/写）
\label{fielddesc:SPIOUTTOTALEOFINTCLR}\item [SPI\_OUT\_TOTAL\_EOF\_INT\_CLR] 所有数据包全部发送到主机后将触发中断，此位为该中断的清除位。可在 CONF 状态下配置。（读/写）
\label{fielddesc:SPIINFIFOFULLERRINTCLR}\item [SPI\_INFIFO\_FULL\_ERR\_INT\_CLR] 1：清除 SPI\_INFIFO\_FULL\_ERR\_INT\_RAW。0：无效。CONF\_buf 可更改该设置。（读/写）
\label{fielddesc:SPIOUTFIFOEMPTYERRINTCLR}\item [SPI\_OUTFIFO\_EMPTY\_ERR\_INT\_CLR] 1：清除 SPI\_OUTFIFO\_EMPTY\_ERR\_INT\_RAW 信号。0：无效。CONF\_buf 可更改该设置。（读/写）
\label{fielddesc:SPISLVCMD6INTCLR}\item [SPI\_SLV\_CMD6\_INT\_CLR] SPI 从机 CMD6 中断清除位。（读/写）
\label{fielddesc:SPISLVCMD7INTCLR}\item [SPI\_SLV\_CMD7\_INT\_CLR] SPI 从机 CMD7 中断清除位。（读/写）
\label{fielddesc:SPISLVCMD8INTCLR}\item [SPI\_SLV\_CMD8\_INT\_CLR] SPI 从机 CMD8 中断清除位。（读/写）
\label{fielddesc:SPISLVCMD9INTCLR}\item [SPI\_SLV\_CMD9\_INT\_CLR] SPI 从机 CMD9 中断清除位。（读/写）
\label{fielddesc:SPISLVCMDAINTCLR}\item [SPI\_SLV\_CMDA\_INT\_CLR] SPI 从机 CMDA 中断清除位。（读/写）
\end{reglist}\end{regdesc}
\end{register}


\begin{register}{H}{SPI\_W0\_REG}{\hexprefix{}0098}\label{regdesc:SPIW0REG}
\regfield{SPI\_BUF0}{32}{0}{{0}}%
\reglabel{Reset}\regnewline%
\begin{regdesc}\begin{reglist}
\label{fielddesc:SPIBUF0}\item [SPI\_BUF0] 数据 buffer 0，32 位，传输单位：字节。从机半双工模式下，可字节寻址。（读/写）
\end{reglist}\end{regdesc}
\end{register}


\begin{register}{H}{SPI\_W1\_REG}{\hexprefix{}009C}\label{regdesc:SPIW1REG}
\regfield{SPI\_BUF1}{32}{0}{{0}}%
\reglabel{Reset}\regnewline%
\begin{regdesc}\begin{reglist}
\label{fielddesc:SPIBUF1}\item [SPI\_BUF1] 数据 buffer 1，32 位，传输单位：字节。从机半双工模式下，可字节寻址。（读/写）
\end{reglist}\end{regdesc}
\end{register}


\begin{register}{H}{SPI\_W2\_REG}{\hexprefix{}00A0}\label{regdesc:SPIW2REG}
\regfield{SPI\_BUF2}{32}{0}{{0}}%
\reglabel{Reset}\regnewline%
\begin{regdesc}\begin{reglist}
\label{fielddesc:SPIBUF2}\item [SPI\_BUF2] 数据 buffer 2，32 位，传输单位：字节。从机半双工模式下，可字节寻址。（读/写）
\end{reglist}\end{regdesc}
\end{register}


\begin{register}{H}{SPI\_W3\_REG}{\hexprefix{}00A4}\label{regdesc:SPIW3REG}
\regfield{SPI\_BUF3}{32}{0}{{0}}%
\reglabel{Reset}\regnewline%
\begin{regdesc}\begin{reglist}
\label{fielddesc:SPIBUF3}\item [SPI\_BUF3] 数据 buffer 3，32 位，传输单位：字节。从机半双工模式下，可字节寻址。（读/写）
\end{reglist}\end{regdesc}
\end{register}


\begin{register}{H}{SPI\_W4\_REG}{\hexprefix{}00A8}\label{regdesc:SPIW4REG}
\regfield{SPI\_BUF4}{32}{0}{{0}}%
\reglabel{Reset}\regnewline%
\begin{regdesc}\begin{reglist}
\label{fielddesc:SPIBUF4}\item [SPI\_BUF4] 数据 buffer 4，32 位，传输单位：字节。从机半双工模式下，可字节寻址。（读/写）
\end{reglist}\end{regdesc}
\end{register}


\begin{register}{H}{SPI\_W5\_REG}{\hexprefix{}00AC}\label{regdesc:SPIW5REG}
\regfield{SPI\_BUF5}{32}{0}{{0}}%
\reglabel{Reset}\regnewline%
\begin{regdesc}\begin{reglist}
\label{fielddesc:SPIBUF5}\item [SPI\_BUF5] 数据 buffer 5，32 位，传输单位：字节。从机半双工模式下，可字节寻址。（读/写）
\end{reglist}\end{regdesc}
\end{register}


\begin{register}{H}{SPI\_W6\_REG}{\hexprefix{}00B0}\label{regdesc:SPIW6REG}
\regfield{SPI\_BUF6}{32}{0}{{0}}%
\reglabel{Reset}\regnewline%
\begin{regdesc}\begin{reglist}
\label{fielddesc:SPIBUF6}\item [SPI\_BUF6] 数据 buffer 6，32 位，传输单位：字节。从机半双工模式下，可字节寻址。（读/写）
\end{reglist}\end{regdesc}
\end{register}


\begin{register}{H}{SPI\_W7\_REG}{\hexprefix{}00B4}\label{regdesc:SPIW7REG}
\regfield{SPI\_BUF7}{32}{0}{{0}}%
\reglabel{Reset}\regnewline%
\begin{regdesc}\begin{reglist}
\label{fielddesc:SPIBUF7}\item [SPI\_BUF7] 数据 buffer 7，32 位，传输单位：字节。从机半双工模式下，可字节寻址。（读/写）
\end{reglist}\end{regdesc}
\end{register}


\begin{register}{H}{SPI\_W8\_REG}{\hexprefix{}00B8}\label{regdesc:SPIW8REG}
\regfield{SPI\_BUF8}{32}{0}{{0}}%
\reglabel{Reset}\regnewline%
\begin{regdesc}\begin{reglist}
\label{fielddesc:SPIBUF8}\item [SPI\_BUF8] 数据 buffer 8，32 位，传输单位：字节。从机半双工模式下，可字节寻址。（读/写）
\end{reglist}\end{regdesc}
\end{register}


\begin{register}{H}{SPI\_W9\_REG}{\hexprefix{}00BC}\label{regdesc:SPIW9REG}
\regfield{SPI\_BUF9}{32}{0}{{0}}%
\reglabel{Reset}\regnewline%
\begin{regdesc}\begin{reglist}
\label{fielddesc:SPIBUF9}\item [SPI\_BUF9] 数据 buffer 9，32 位，传输单位：字节。从机半双工模式下，可字节寻址。（读/写）
\end{reglist}\end{regdesc}
\end{register}


\begin{register}{H}{SPI\_W10\_REG}{\hexprefix{}00C0}\label{regdesc:SPIW10REG}
\regfield{SPI\_BUF10}{32}{0}{{0}}%
\reglabel{Reset}\regnewline%
\begin{regdesc}\begin{reglist}
\label{fielddesc:SPIBUF10}\item [SPI\_BUF10] 数据 buffer 10，32 位，传输单位：字节。从机半双工模式下，可字节寻址。（读/写）
\end{reglist}\end{regdesc}
\end{register}


\begin{register}{H}{SPI\_W11\_REG}{\hexprefix{}00C4}\label{regdesc:SPIW11REG}
\regfield{SPI\_BUF11}{32}{0}{{0}}%
\reglabel{Reset}\regnewline%
\begin{regdesc}\begin{reglist}
\label{fielddesc:SPIBUF11}\item [SPI\_BUF11] 数据 buffer 11，32 位，传输单位：字节。从机半双工模式下，可字节寻址。（读/写）
\end{reglist}\end{regdesc}
\end{register}


\begin{register}{H}{SPI\_W12\_REG}{\hexprefix{}00C8}\label{regdesc:SPIW12REG}
\regfield{SPI\_BUF12}{32}{0}{{0}}%
\reglabel{Reset}\regnewline%
\begin{regdesc}\begin{reglist}
\label{fielddesc:SPIBUF12}\item [SPI\_BUF12] 数据 buffer 12，32 位，传输单位：字节。从机半双工模式下，可字节寻址。（读/写）
\end{reglist}\end{regdesc}
\end{register}


\begin{register}{H}{SPI\_W13\_REG}{\hexprefix{}00CC}\label{regdesc:SPIW13REG}
\regfield{SPI\_BUF13}{32}{0}{{0}}%
\reglabel{Reset}\regnewline%
\begin{regdesc}\begin{reglist}
\label{fielddesc:SPIBUF13}\item [SPI\_BUF13] 数据 buffer 13，32 位，传输单位：字节。从机半双工模式下，可字节寻址。（读/写）
\end{reglist}\end{regdesc}
\end{register}


\begin{register}{H}{SPI\_W14\_REG}{\hexprefix{}00D0}\label{regdesc:SPIW14REG}
\regfield{SPI\_BUF14}{32}{0}{{0}}%
\reglabel{Reset}\regnewline%
\begin{regdesc}\begin{reglist}
\label{fielddesc:SPIBUF14}\item [SPI\_BUF14] 数据 buffer 14，32 位，传输单位：字节。从机半双工模式下，可字节寻址。（读/写）
\end{reglist}\end{regdesc}
\end{register}


\begin{register}{H}{SPI\_W15\_REG}{\hexprefix{}00D4}\label{regdesc:SPIW15REG}
\regfield{SPI\_BUF15}{32}{0}{{0}}%
\reglabel{Reset}\regnewline%
\begin{regdesc}\begin{reglist}
\label{fielddesc:SPIBUF15}\item [SPI\_BUF15] 数据 buffer 15，32 位，传输单位：字节。从机半双工模式下，可字节寻址。（读/写）
\end{reglist}\end{regdesc}
\end{register}


\begin{register}{H}{SPI\_W16\_REG}{\hexprefix{}00D8}\label{regdesc:SPIW16REG}
\regfield{SPI\_BUF16}{32}{0}{{0}}%
\reglabel{Reset}\regnewline%
\begin{regdesc}\begin{reglist}
\label{fielddesc:SPIBUF16}\item [SPI\_BUF16] 数据 buffer 16，32 位，传输单位：字节。从机半双工模式下，可字节寻址。（读/写）
\end{reglist}\end{regdesc}
\end{register}


\begin{register}{H}{SPI\_W17\_REG}{\hexprefix{}00DC}\label{regdesc:SPIW17REG}
\regfield{SPI\_BUF17}{32}{0}{{0}}%
\reglabel{Reset}\regnewline%
\begin{regdesc}\begin{reglist}
\label{fielddesc:SPIBUF17}\item [SPI\_BUF17] 数据 buffer 17，32 位，传输单位：字节。从机半双工模式下，可字节寻址。（读/写）
\end{reglist}\end{regdesc}
\end{register}


\begin{register}{H}{SPI\_DIN\_MODE\_REG}{\hexprefix{}00E0}\label{regdesc:SPIDINMODEREG}
\regfield{(reserved)}{7}{25}{0000000}%
\regfield{SPI\_TIMING\_CLK\_ENA}{1}{24}{{0}}%
\regfield{SPI\_DIN7\_MODE}{3}{21}{{\hexprefix{}0}}%
\regfield{SPI\_DIN6\_MODE}{3}{18}{{\hexprefix{}0}}%
\regfield{SPI\_DIN5\_MODE}{3}{15}{{\hexprefix{}0}}%
\regfield{SPI\_DIN4\_MODE}{3}{12}{{\hexprefix{}0}}%
\regfield{SPI\_DIN3\_MODE}{3}{9}{{\hexprefix{}0}}%
\regfield{SPI\_DIN2\_MODE}{3}{6}{{\hexprefix{}0}}%
\regfield{SPI\_DIN1\_MODE}{3}{3}{{\hexprefix{}0}}%
\regfield{SPI\_DIN0\_MODE}{3}{0}{{\hexprefix{}0}}%
\reglabel{Reset}\regnewline%
\begin{regdesc}\begin{reglist}
\label{fielddesc:SPIDIN0MODE}\item [SPI\_DIN0\_MODE] 选择不同的时钟源，对主机输入信号 FSPID 进行延迟后才采样。0：无输入延迟；1：APB\_CLK 上升沿；2：APB\_CLK 下降沿；3：HCLK 上升沿；4: HCLK 下降沿。可在 CONF 阶段配置。（读/写）
\label{fielddesc:SPIDIN1MODE}\item [SPI\_DIN1\_MODE] 选择不同的时钟源，对主机输入信号 FSPIQ 进行延迟后才采样。0：无输入延迟；1：APB\_CLK 上升沿；2：APB\_CLK 下降沿；3：HCLK 上升沿；4: HCLK 下降沿。可在 CONF 阶段配置。（读/写）
\label{fielddesc:SPIDIN2MODE}\item [SPI\_DIN2\_MODE] 选择不同的时钟源，对主机输入信号 FSPIWP 进行延迟后才采样。0：无输入延迟；1：APB\_CLK 上升沿；2：APB\_CLK 下降沿；3：HCLK 上升沿；4: HCLK 下降沿。可在 CONF 阶段配置。（读/写）
\label{fielddesc:SPIDIN3MODE}\item [SPI\_DIN3\_MODE] 选择不同的时钟源，对主机输入信号 FSPIHD 进行延迟后才采样。0：无输入延迟；1：APB\_CLK 上升沿；2：APB\_CLK 下降沿；3：HCLK 上升沿；4: HCLK 下降沿。可在 CONF 阶段配置。（读/写）
\label{fielddesc:SPIDIN4MODE}\item [SPI\_DIN4\_MODE] 选择不同的时钟源，对主机输入信号 FSPIIO4 进行延迟后才采样。0：无输入延迟；1：APB\_CLK 上升沿；2：APB\_CLK 下降沿；3：HCLK 上升沿；4: HCLK 下降沿。可在 CONF 阶段配置。（读/写）
\label{fielddesc:SPIDIN5MODE}\item [SPI\_DIN5\_MODE] 选择不同的时钟源，对主机输入信号 FSPIIO5 进行延迟后才采样。0：无输入延迟；1：APB\_CLK 上升沿；2：APB\_CLK 下降沿；3：HCLK 上升沿；4: HCLK 下降沿。可在 CONF 阶段配置。（读/写）
\label{fielddesc:SPIDIN6MODE}\item [SPI\_DIN6\_MODE] 选择不同的时钟源，对主机输入信号 FSPIIO6 进行延迟后才采样。0：无输入延迟；1：APB\_CLK 上升沿；2：APB\_CLK 下降沿；3：HCLK 上升沿；4: HCLK 下降沿。可在 CONF 阶段配置。（读/写）
\label{fielddesc:SPIDIN7MODE}\item [SPI\_DIN7\_MODE] 选择不同的时钟源，对主机输入信号 FSPIIO7 进行延迟后才采样。0：无输入延迟；1：APB\_CLK 上升沿；2：APB\_CLK 下降沿；3：HCLK 上升沿；4: HCLK 下降沿。可在 CONF 阶段配置。（读/写）
\label{fielddesc:SPITIMINGCLKENA}\item [SPI\_TIMING\_CLK\_ENA] 1：使能时序模块的高频时钟 HCLK 160 MHz；0：禁用该功能。可在 CONF 阶段配置。（读/写）
\end{reglist}\end{regdesc}
\end{register}


\begin{register}{H}{SPI\_DIN\_NUM\_REG}{\hexprefix{}00E4}\label{regdesc:SPIDINNUMREG}
\regfield{(reserved)}{16}{16}{0000000000000000}%
\regfield{SPI\_DIN7\_NUM}{2}{14}{{\hexprefix{}0}}%
\regfield{SPI\_DIN6\_NUM}{2}{12}{{\hexprefix{}0}}%
\regfield{SPI\_DIN5\_NUM}{2}{10}{{\hexprefix{}0}}%
\regfield{SPI\_DIN4\_NUM}{2}{8}{{\hexprefix{}0}}%
\regfield{SPI\_DIN3\_NUM}{2}{6}{{\hexprefix{}0}}%
\regfield{SPI\_DIN2\_NUM}{2}{4}{{\hexprefix{}0}}%
\regfield{SPI\_DIN1\_NUM}{2}{2}{{\hexprefix{}0}}%
\regfield{SPI\_DIN0\_NUM}{2}{0}{{\hexprefix{}0}}%
\reglabel{Reset}\regnewline%
\begin{regdesc}\begin{reglist}
\label{fielddesc:SPIDIN0NUM}\item [SPI\_DIN0\_NUM] 针对 SPI\_DIN0\_MODE 选定的时钟源，配置对输入信号 FSPID 的延迟周期数。0：延迟 1 个周期；1：延迟 2 个周期，以此类推。可在 CONF 状态下配置。（读/写）
\label{fielddesc:SPIDIN1NUM}\item [SPI\_DIN1\_NUM] 针对 SPI\_DIN1\_MODE 选定的时钟源，配置对输入信号 FSPIQ 的延迟周期数。0：延迟 1 个周期；1：延迟 2 个周期，以此类推。可在 CONF 状态下配置。（读/写）
\label{fielddesc:SPIDIN2NUM}\item [SPI\_DIN2\_NUM] 针对 SPI\_DIN2\_MODE 选定的时钟源，配置对输入信号 FSPIWP 的延迟周期数。0：延迟 1 个周期；1：延迟 2 个周期，以此类推。可在 CONF 状态下配置。（读/写）
\label{fielddesc:SPIDIN3NUM}\item [SPI\_DIN3\_NUM] 针对 SPI\_DIN3\_MODE 选定的时钟源，配置对输入信号 FSPIHD 的延迟周期数。0：延迟 1 个周期；1：延迟 2 个周期，以此类推。可在 CONF 状态下配置。（读/写）
\label{fielddesc:SPIDIN4NUM}\item [SPI\_DIN4\_NUM] 针对 SPI\_DIN4\_MODE 选定的时钟源，配置对输入信号 FSPIIO4 的延迟周期数。0：延迟 1 个周期；1：延迟 2 个周期，以此类推。可在 CONF 状态下配置。（读/写）
\label{fielddesc:SPIDIN5NUM}\item [SPI\_DIN5\_NUM] 针对 SPI\_DIN5\_MODE 选定的时钟源，配置对输入信号 FSPIIO5 的延迟周期数。0：延迟 1 个周期；1：延迟 2 个周期，以此类推。可在 CONF 状态下配置。（读/写）
\label{fielddesc:SPIDIN6NUM}\item [SPI\_DIN6\_NUM] 针对 SPI\_DIN6\_MODE 选定的时钟源，配置对输入信号 FSPIIO6 的延迟周期数。0：延迟 1 个周期；1：延迟 2 个周期，以此类推。可在 CONF 状态下配置。（读/写）
\label{fielddesc:SPIDIN7NUM}\item [SPI\_DIN7\_NUM] 针对 SPI\_DIN7\_MODE 选定的时钟源，配置对输入信号 FSPIIO7 的延迟周期数。0：延迟 1 个周期；1：延迟 2 个周期，以此类推。可在 CONF 状态下配置。（读/写）
\end{reglist}\end{regdesc}
\end{register}


\begin{register}{H}{SPI\_DOUT\_MODE\_REG}{\hexprefix{}00E8}\label{regdesc:SPIDOUTMODEREG}
\regfield{(reserved)}{8}{24}{00000000}%
\regfield{SPI\_DOUT7\_MODE}{3}{21}{{\hexprefix{}0}}%
\regfield{SPI\_DOUT6\_MODE}{3}{18}{{\hexprefix{}0}}%
\regfield{SPI\_DOUT5\_MODE}{3}{15}{{\hexprefix{}0}}%
\regfield{SPI\_DOUT4\_MODE}{3}{12}{{\hexprefix{}0}}%
\regfield{SPI\_DOUT3\_MODE}{3}{9}{{\hexprefix{}0}}%
\regfield{SPI\_DOUT2\_MODE}{3}{6}{{\hexprefix{}0}}%
\regfield{SPI\_DOUT1\_MODE}{3}{3}{{\hexprefix{}0}}%
\regfield{SPI\_DOUT0\_MODE}{3}{0}{{\hexprefix{}0}}%
\reglabel{Reset}\regnewline%
\begin{regdesc}\begin{reglist}
\label{fielddesc:SPIDOUT0MODE}\item [SPI\_DOUT0\_MODE] 选择不同的时钟源，对主机输出信号 FSPID 进行延迟后才输出。0：无输出延迟；1：APB\_CLK 上升沿；2：APB\_CLK 下降沿；3：HCLK 上升沿；4: HCLK 下降沿。可在 CONF 阶段配置。（读/写）
\label{fielddesc:SPIDOUT1MODE}\item [SPI\_DOUT1\_MODE] 选择不同的时钟源，对主机输出信号 FSPIQ 进行延迟后才输出。0：无输出延迟；1：APB\_CLK 上升沿；2：APB\_CLK 下降沿；3：HCLK 上升沿；4: HCLK 下降沿。可在 CONF 阶段配置。（读/写）
\label{fielddesc:SPIDOUT2MODE}\item [SPI\_DOUT2\_MODE] 选择不同的时钟源，对主机输出信号 FSPIWP 进行延迟后才输出。0：无输出延迟；1：APB\_CLK 上升沿；2：APB\_CLK 下降沿；3：HCLK 上升沿；4: HCLK 下降沿。可在 CONF 阶段配置。（读/写）
\label{fielddesc:SPIDOUT3MODE}\item [SPI\_DOUT3\_MODE] 选择不同的时钟源，对主机输出信号 FSPIHD 进行延迟后才输出。0：无输出延迟；1：APB\_CLK 上升沿；2：APB\_CLK 下降沿；3：HCLK 上升沿；4: HCLK 下降沿。可在 CONF 阶段配置。（读/写）
\label{fielddesc:SPIDOUT4MODE}\item [SPI\_DOUT4\_MODE] 选择不同的时钟源，对主机输出信号 FSPIIO4 进行延迟后才输出。0：无输出延迟；1：APB\_CLK 上升沿；2：APB\_CLK 下降沿；3：HCLK 上升沿；4: HCLK 下降沿。可在 CONF 阶段配置。（读/写）
\label{fielddesc:SPIDOUT5MODE}\item [SPI\_DOUT5\_MODE] 选择不同的时钟源，对主机输出信号 FSPIIO5 进行延迟后才输出。0：无输出延迟；1：APB\_CLK 上升沿；2：APB\_CLK 下降沿；3：HCLK 上升沿；4: HCLK 下降沿。可在 CONF 阶段配置。（读/写）
\label{fielddesc:SPIDOUT6MODE}\item [SPI\_DOUT6\_MODE] 选择不同的时钟源，对主机输出信号 FSPIIO6 进行延迟后才输出。0：无输出延迟；1：APB\_CLK 上升沿；2：APB\_CLK 下降沿；3：HCLK 上升沿；4: HCLK 下降沿。可在 CONF 阶段配置。（读/写）
\label{fielddesc:SPIDOUT7MODE}\item [SPI\_DOUT7\_MODE] 选择不同的时钟源，对主机输出信号 FSPIIO7 进行延迟后才输出。0：无输出延迟；1：APB\_CLK 上升沿；2：APB\_CLK 下降沿；3：HCLK 上升沿；4: HCLK 下降沿。可在 CONF 阶段配置。（读/写）
\end{reglist}\end{regdesc}
\end{register}


\begin{register}{H}{SPI\_DOUT\_NUM\_REG}{\hexprefix{}00EC}\label{regdesc:SPIDOUTNUMREG}
\regfield{(reserved)}{16}{16}{0000000000000000}%
\regfield{SPI\_DOUT7\_NUM}{2}{14}{{\hexprefix{}0}}%
\regfield{SPI\_DOUT6\_NUM}{2}{12}{{\hexprefix{}0}}%
\regfield{SPI\_DOUT5\_NUM}{2}{10}{{\hexprefix{}0}}%
\regfield{SPI\_DOUT4\_NUM}{2}{8}{{\hexprefix{}0}}%
\regfield{SPI\_DOUT3\_NUM}{2}{6}{{\hexprefix{}0}}%
\regfield{SPI\_DOUT2\_NUM}{2}{4}{{\hexprefix{}0}}%
\regfield{SPI\_DOUT1\_NUM}{2}{2}{{\hexprefix{}0}}%
\regfield{SPI\_DOUT0\_NUM}{2}{0}{{\hexprefix{}0}}%
\reglabel{Reset}\regnewline%
\begin{regdesc}\begin{reglist}
\label{fielddesc:SPIDOUT0NUM}\item [SPI\_DOUT0\_NUM] 针对 SPI\_DOUT0\_MODE 选定的时钟源，配置对输出信号 FSPID 的延迟周期数。0：延迟 1 个周期；1：延迟 2 个周期，以此类推。可在 CONF 状态下配置。（读/写）
\label{fielddesc:SPIDOUT1NUM}\item [SPI\_DOUT1\_NUM] 针对 SPI\_DOUT1\_MODE 选定的时钟源，配置对输出信号 FSPIQ 的延迟周期数。0：延迟 1 个周期；1：延迟 2 个周期，以此类推。可在 CONF 状态下配置。（读/写）
\label{fielddesc:SPIDOUT2NUM}\item [SPI\_DOUT2\_NUM] 针对 SPI\_DOUT2\_MODE 选定的时钟源，配置对输出信号 FSPIWP 的延迟周期数。0：延迟 1 个周期；1：延迟 2 个周期，以此类推。可在 CONF 状态下配置。（读/写）
\label{fielddesc:SPIDOUT3NUM}\item [SPI\_DOUT3\_NUM] 针对 SPI\_DOUT3\_MODE 选定的时钟源，配置对输出信号 FSPIHD 的延迟周期数。0：延迟 1 个周期；1：延迟 2 个周期，以此类推。可在 CONF 状态下配置。（读/写）
\label{fielddesc:SPIDOUT4NUM}\item [SPI\_DOUT4\_NUM] 针对 SPI\_DOUT4\_MODE 选定的时钟源，配置对输出信号 FSPIIO4 的延迟周期数。0：延迟 1 个周期；1：延迟 2 个周期，以此类推。可在 CONF 状态下配置。（读/写）
\label{fielddesc:SPIDOUT5NUM}\item [SPI\_DOUT5\_NUM] 针对 SPI\_DOUT5\_MODE 选定的时钟源，配置对输出信号 FSPIIO5 的延迟周期数。0：延迟 1 个周期；1：延迟 2 个周期，以此类推。可在 CONF 状态下配置。（读/写）
\label{fielddesc:SPIDOUT6NUM}\item [SPI\_DOUT6\_NUM] 针对 SPI\_DOUT6\_MODE 选定的时钟源，配置对输出信号 FSPIIO6 的延迟周期数。0：延迟 1 个周期；1：延迟 2 个周期，以此类推。可在 CONF 状态下配置。（读/写）
\label{fielddesc:SPIDOUT7NUM}\item [SPI\_DOUT7\_NUM] 针对 SPI\_DOUT7\_MODE 选定的时钟源，配置对输出信号 FSPIIO7 的延迟周期数。0：延迟 1 个周期；1：延迟 2 个周期，以此类推。可在 CONF 状态下配置。（读/写）
\end{reglist}\end{regdesc}
\end{register}


\begin{register}{H}{SPI\_LCD\_CTRL\_REG}{\hexprefix{}00F0}\label{regdesc:SPILCDCTRLREG}
\regfield{SPI\_LCD\_MODE\_EN}{1}{31}{{0}}%
\regfield{SPI\_LCD\_VT\_HEIGHT}{10}{21}{{0}}%
\regfield{SPI\_LCD\_VA\_HEIGHT}{10}{11}{{0}}%
\regfield{SPI\_LCD\_HB\_FRONT}{11}{0}{{0}}%
\reglabel{Reset}\regnewline%
\begin{regdesc}\begin{reglist}
\label{fielddesc:SPILCDHBFRONT}\item [SPI\_LCD\_HB\_FRONT] 配置 (\hyperref[fielddesc:SPILCDHSYNCPOSITION]{HSYNC\_POSITION} + \hyperref[fielddesc:SPILCDHSYNCWIDTH]{HSYNC\_WIDTH} + 水平后肩) 的值。可在 CONF 阶段配置。（读/写）
\label{fielddesc:SPILCDVAHEIGHT}\item [SPI\_LCD\_VA\_HEIGHT] 帧有效垂直高度。可在 CONF 阶段配置。（读/写）
\label{fielddesc:SPILCDVTHEIGHT}\item [SPI\_LCD\_VT\_HEIGHT] 帧总垂直高度。可在 CONF 阶段配置。（读/写）
\label{fielddesc:SPILCDMODEEN}\item [SPI\_LCD\_MODE\_EN] 1：使能 LCD 模式，输出 VSYNC、HSYNC 和 DE。0：禁用。可在 CONF 阶段配置。（读/写）
\end{reglist}\end{regdesc}
\end{register}


\begin{register}{H}{SPI\_LCD\_CTRL1\_REG}{\hexprefix{}00F4}\label{regdesc:SPILCDCTRL1REG}
\regfield{SPI\_LCD\_HT\_WIDTH}{12}{20}{{0}}%
\regfield{SPI\_LCD\_HA\_WIDTH}{12}{8}{{0}}%
\regfield{SPI\_LCD\_VB\_FRONT}{8}{0}{{0}}%
\reglabel{Reset}\regnewline%
\begin{regdesc}\begin{reglist}
\label{fielddesc:SPILCDVBFRONT}\item [SPI\_LCD\_VB\_FRONT] 配置 (\hyperref[fielddesc:SPILCDVSYNCWIDTH]{VSYNC\_WIDTH} + 垂直后肩) 的值。可在 CONF 阶段配置。（读/写）
\label{fielddesc:SPILCDHAWIDTH}\item [SPI\_LCD\_HA\_WIDTH] 帧有效水平宽度。可在 CONF 阶段配置。（读/写）
\label{fielddesc:SPILCDHTWIDTH}\item [SPI\_LCD\_HT\_WIDTH] 帧总水平宽度。可在 CONF 阶段配置。（读/写）
\end{reglist}\end{regdesc}
\end{register}


\begin{register}{H}{SPI\_LCD\_CTRL2\_REG}{\hexprefix{}00F8}\label{regdesc:SPILCDCTRL2REG}
\regfield{SPI\_LCD\_HSYNC\_POSITION}{8}{24}{{0}}%
\regfield{SPI\_HSYNC\_IDLE\_POL}{1}{23}{{0}}%
\regfield{SPI\_LCD\_HSYNC\_WIDTH}{7}{16}{{1}}%
\regfield{(reserved)}{8}{8}{00000000}%
\regfield{SPI\_VSYNC\_IDLE\_POL}{1}{7}{{0}}%
\regfield{SPI\_LCD\_VSYNC\_WIDTH}{7}{0}{{1}}%
\reglabel{Reset}\regnewline%
\begin{regdesc}\begin{reglist}
\label{fielddesc:SPILCDVSYNCWIDTH}\item [SPI\_LCD\_VSYNC\_WIDTH] 信号线上 SPI\_VSYNC 有效脉冲宽度。可在 CONF 阶段配置。（读/写）
\label{fielddesc:SPIVSYNCIDLEPOL}\item [SPI\_VSYNC\_IDLE\_POL] SPI\_VSYNC 空闲值。可在 CONF 阶段配置。（读/写）
\label{fielddesc:SPILCDHSYNCWIDTH}\item [SPI\_LCD\_HSYNC\_WIDTH] 信号线上 SPI\_HSYNC 有效脉冲宽度。可在 CONF 阶段配置。参数值=寄存器值 + 1。（读/写）
\label{fielddesc:SPIHSYNCIDLEPOL}\item [SPI\_HSYNC\_IDLE\_POL] SPI\_HSYNC 空闲值。可在 CONF 阶段配置。（读/写）
\label{fielddesc:SPILCDHSYNCPOSITION}\item [SPI\_LCD\_HSYNC\_POSITION] 信号线上 SPI\_HSYNC 有效脉冲位置。可在 CONF 状态下配置。参数值=寄存器值 + 1。单位为像素。（读/写）
\end{reglist}\end{regdesc}
\end{register}


\begin{register}{H}{SPI\_LCD\_D\_MODE\_REG}{\hexprefix{}00FC}\label{regdesc:SPILCDDMODEREG}
\regfield{(reserved)}{15}{17}{000000000000000}%
\regfield{SPI\_HS\_BLANK\_EN}{1}{16}{{0}}%
\regfield{SPI\_DE\_IDLE\_POL}{1}{15}{{0}}%
\regfield{SPI\_D\_VSYNC\_MODE}{3}{12}{{\hexprefix{}0}}%
\regfield{SPI\_D\_HSYNC\_MODE}{3}{9}{{\hexprefix{}0}}%
\regfield{SPI\_D\_DE\_MODE}{3}{6}{{\hexprefix{}0}}%
\regfield{SPI\_D\_CD\_MODE}{3}{3}{{\hexprefix{}0}}%
\regfield{SPI\_D\_DQS\_MODE}{3}{0}{{\hexprefix{}0}}%
\reglabel{Reset}\regnewline%
\begin{regdesc}\begin{reglist}
\label{fielddesc:SPIDDQSMODE}\item [SPI\_D\_DQS\_MODE] 选择不同的时钟源，对主机输出信号 FSPIDQS 进行延迟后才输出。0：无输出延迟；1：APB\_CLK 上升沿；2：APB\_CLK 下降沿；3：HCLK 上升沿；4: HCLK 下降沿。可在 CONF 阶段配置。（读/写）
\label{fielddesc:SPIDCDMODE}\item [SPI\_D\_CD\_MODE] 选择不同的时钟源，对主机输出信号 FSPICD 进行延迟后才输出。0：无输出延迟；1：APB\_CLK 上升沿；2：APB\_CLK 下降沿；3：HCLK 上升沿；4: HCLK 下降沿。可在 CONF 阶段配置。（读/写）
\label{fielddesc:SPIDDEMODE}\item [SPI\_D\_DE\_MODE] 选择不同的时钟源，对主机输出信号 FSPI\_DE 进行延迟后才输出。0：无输出延迟；1：APB\_CLK 上升沿；2：APB\_CLK 下降沿；3：HCLK 上升沿；4: HCLK 下降沿。可在 CONF 阶段配置。（读/写）
\label{fielddesc:SPIDHSYNCMODE}\item [SPI\_D\_HSYNC\_MODE] 选择不同的时钟源，对主机输出信号 FSPI\_HSYNC 进行延迟后才输出。0：无输出延迟；1：APB\_CLK 上升沿；2：APB\_CLK 下降沿；3：HCLK 上升沿；4: HCLK 下降沿。可在 CONF 阶段配置。（读/写）
\label{fielddesc:SPIDVSYNCMODE}\item [SPI\_D\_VSYNC\_MODE] 选择不同的时钟源，对主机输出信号 FSPI\_VSYNC 进行延迟后才输出。0：无输出延迟；1：APB\_CLK 上升沿；2：APB\_CLK 下降沿；3：HCLK 上升沿；4: HCLK 下降沿。可在 CONF 阶段配置。（读/写）
\label{fielddesc:SPIDEIDLEPOL}\item [SPI\_DE\_IDLE\_POL] SPI\_DE 空闲值。（读/写）
\label{fielddesc:SPIHSBLANKEN}\item [SPI\_HS\_BLANK\_EN] 1：在分段配置传输过程中，或单次传输过程中，在垂直消隐行输出 SPI\_HSYNC。0：SPI\_HSYNC 脉冲仅在分段配置传输模式下活动区域中的行有效。（读/写）
\end{reglist}\end{regdesc}
\end{register}


\begin{register}{H}{SPI\_LCD\_D\_NUM\_REG}{\hexprefix{}0100}\label{regdesc:SPILCDDNUMREG}
\regfield{(reserved)}{22}{10}{0000000000000000000000}%
\regfield{SPI\_D\_VSYNC\_NUM}{2}{8}{{\hexprefix{}0}}%
\regfield{SPI\_D\_HSYNC\_NUM}{2}{6}{{\hexprefix{}0}}%
\regfield{SPI\_D\_DE\_NUM}{2}{4}{{\hexprefix{}0}}%
\regfield{SPI\_D\_CD\_NUM}{2}{2}{{\hexprefix{}0}}%
\regfield{SPI\_D\_DQS\_NUM}{2}{0}{{\hexprefix{}0}}%
\reglabel{Reset}\regnewline%
\begin{regdesc}\begin{reglist}
\label{fielddesc:SPIDDQSNUM}\item [SPI\_D\_DQS\_NUM] 针对 SPI\_D\_DQS\_MODE 选定的时钟源，配置对输出信号 FSPIDQS 的延迟周期数。0：延迟 1 个周期；1：延迟 2 个周期，以此类推。可在 CONF 阶段配置。（读/写）
\label{fielddesc:SPIDCDNUM}\item [SPI\_D\_CD\_NUM] 针对 SPI\_D\_CD\_MODE 选定的时钟源，配置对输出信号 FSPI\_CD 的延迟周期数。0：延迟 1 个周期；1：延迟 2 个周期，以此类推。可在 CONF 阶段配置。（读/写）
\label{fielddesc:SPIDDENUM}\item [SPI\_D\_DE\_NUM] 针对 SPI\_D\_DE\_MODE 选定的时钟源，配置对输出信号 FSPI\_DE 的延迟周期数。0：延迟 1 个周期；1：延迟 2 个周期，以此类推。可在 CONF 阶段配置。（读/写）
\label{fielddesc:SPIDHSYNCNUM}\item [SPI\_D\_HSYNC\_NUM] 针对 SPI\_D\_HSYNC\_MODE 选定的时钟源，配置对输出信号 FSPI\_HSYNC 的延迟周期数。0：延迟 1 个周期；1：延迟 2 个周期，以此类推。可在 CONF 阶段配置。（读/写）
\label{fielddesc:SPIDVSYNCNUM}\item [SPI\_D\_VSYNC\_NUM] 针对 SPI\_D\_VSYNC\_MODE 选定的时钟源，配置对输出信号 FSPI\_VSYNC 的延迟周期数。0：延迟 1 个周期；1：延迟 2 个周期，以此类推。可在 CONF 阶段配置。（读/写）
\end{reglist}\end{regdesc}
\end{register}


\begin{register}{H}{SPI\_DATE\_REG}{\hexprefix{}03FC}\label{regdesc:SPIDATEREG}
\regfield{(reserved)}{4}{28}{0000}%
\regfield{SPI\_DATE}{28}{0}{{\hexprefix{}1907240}}%
\reglabel{Reset}\regnewline%
\begin{regdesc}\begin{reglist}
\label{fielddesc:SPIDATE}\item [SPI\_DATE] 版本控制寄存器。（读/写）
\end{reglist}\end{regdesc}
\end{register}
